\chapter[Równania Laplace'a i Poissona]{Równania Laplace'a i Poissona: funkcje harmoniczne i potencjał}
\label{chap:pde-laplace}

W tym rozdziale badamy dwa równania na zbiorze otwartym $\Omega\subset\R^n$:
$\Delta u=0$ oraz $-\Delta u=f$. Pierwsze opisuje między innymi
stacjonarną temperaturę i potencjały bez źródeł; drugie uwzględnia
źródła. W przeciwieństwie do równania transportu nie wystarcza
znajomość danych na dowolnej hiperpowierzchni. Naturalnym zadaniem
jest zwykle problem Dirichleta na \emph{całym} brzegu.

W całym rozdziale, o ile nie zaznaczono inaczej,
$n\geq2$ jest liczbą całkowitą, a funkcje skalarne są
rzeczywiste. Przez $\lambda_k$ oznaczamy miarę Lebesgue'a
w $\R^k$ dla każdej liczby całkowitej $k\geq1$, przez
$\sigma$ miarę powierzchniową na rozpatrywanej gładkiej
hiperpowierzchni, a przez $\omega_n=\sigma(S^{n-1})$
pole jednostkowej sfery. Ze wzoru na całkowanie
we współrzędnych biegunowych wynika
$\lambda_n(B_r(0))=\omega_n\int_0^r s^{n-1}\,
\mathrm d\lambda_1(s)=\omega_n r^n/n$.
Znak laplasjanu ustalamy jako $\Delta=\sum_{j=1}^n\partial_{x_j}^2$;
rozwiązanie fundamentalne będziemy normalizować przez $-\Delta\Phi=\delta_0$.

\section{Własność wartości średniej}
\label{sec:pde3-mean}

\begin{definicja}[Funkcja harmoniczna]
Niech $\Omega\subset\R^n$ będzie niepustym zbiorem otwartym.
Funkcję rzeczywistą $u\in C^2(\Omega)$ nazywamy \emph{harmoniczną}
na $\Omega$, jeśli $\Delta u=0$ w każdym punkcie $\Omega$.
Funkcję $u\in C^2(\Omega)$ nazywamy \emph{subharmoniczną},
gdy $\Delta u\geq0$, oraz \emph{superharmoniczną},
gdy $\Delta u\leq0$. Później rozszerzymy te pojęcia
na funkcje bez klasycznych drugich pochodnych.
\end{definicja}

Dla $u\in C(\Omega)$, $x_0\in\Omega$ oraz
$0<r<\operatorname{dist}(x_0,\partial\Omega)$ zapisujemy
średnią sferyczną
\begin{equation}\label{eq:pde3-sphere-mean}
 M_u(x_0,r)=\frac{1}{\omega_n r^{n-1}}
 \int_{\partial B_r(x_0)}u(y)\,\mathrm d\sigma(y)
 =\frac1{\omega_n}\int_{S^{n-1}}u(x_0+r\theta)\,
 \mathrm d\sigma(\theta).
\end{equation}
Druga równość wynika ze zmiany zmiennych
$y=x_0+r\theta$: miara powierzchniowa skaluje się
przez czynnik $r^{n-1}$, gdyż pochodna tego
odwzorowania na każdej przestrzeni stycznej
jest mnożeniem przez $r$.
Jeśli $\Omega=\R^n$, przyjmujemy
$\operatorname{dist}(x_0,\partial\Omega)=+\infty$.

\begin{twierdzenie}[Własność wartości średniej]
\label{thm:pde3-mean}
Niech $n\geq2$, $\Omega\subset\R^n$ będzie otwarta,
$u\in C^2(\Omega)$ i $\Delta u=0$ w $\Omega$.
Dla każdego $x_0\in\Omega$ oraz $r>0$ takich, że
$\overline{B_r(x_0)}\subset\Omega$, zachodzi
\begin{equation}\label{eq:pde3-both-means}
 u(x_0)=\frac{1}{\omega_n r^{n-1}}
 \int_{\partial B_r(x_0)}u(y)\,\mathrm d\sigma(y)
 =\frac{1}{\lambda_n(B_r(x_0))}
 \int_{B_r(x_0)}u(y)\,\mathrm d\lambda_n(y).
\end{equation}
\end{twierdzenie}
\begin{proof}
Ustalmy $x_0\in\Omega$ i oznaczmy
$d_0=\operatorname{dist}(x_0,\partial\Omega)$.
Rozważmy $M(r)=M_u(x_0,r)$ dla $0<r<d_0$.
Dla każdego zwartego przedziału $[a,b]\subset(0,d_0)$
funkcja $(r,\theta)\mapsto u(x_0+r\theta)$ oraz jej
pochodna $\nabla u(x_0+r\theta)\cdot\theta$ są ciągłe
na $[a,b]\times S^{n-1}$, a więc ograniczone.
Ponieważ sfera ma skończoną miarę, można różniczkować
pod znakiem całki. Na $\partial B_r(x_0)$ jednostkowa
normalna zewnętrzna wynosi $\nu(y)=(y-x_0)/r$.
Zmiana zmiennych $y=x_0+r\theta$, a następnie
twierdzenie Gaussa zastosowane do pola $\nabla u$
na kuli dają
\[
 M'(r)=\frac{1}{\omega_n r^{n-1}}
 \int_{\partial B_r(x_0)}\partial_\nu u(y)\,
 \mathrm d\sigma(y)
 =\frac{1}{\omega_n r^{n-1}}
 \int_{B_r(x_0)}\Delta u(y)\,\mathrm d\lambda_n(y)=0.
\]
Zatem $M$ jest stała na $(0,d_0)$. Ponadto
\[
 |M(r)-u(x_0)|\leq
 \sup_{\theta\in S^{n-1}}|u(x_0+r\theta)-u(x_0)|
 \longrightarrow0\quad(r\downarrow0)
\]
z ciągłości $u$ w $x_0$, co daje średnią sferyczną.
Dla $\overline{B_r(x_0)}\subset\Omega$ mamy $r<d_0$.
We współrzędnych biegunowych otrzymujemy
\[
 \begin{aligned}
 \int_{B_r(x_0)}u(y)\,\mathrm d\lambda_n(y)
 &=\omega_n\int_0^r M(s)s^{n-1}\,\mathrm d\lambda_1(s)\\
 &=u(x_0)\frac{\omega_n r^n}{n}
 =u(x_0)\lambda_n(B_r(x_0)).
 \end{aligned}
\]
Po podzieleniu przez dodatnią miarę kuli otrzymujemy
drugą równość w~\eqref{eq:pde3-both-means}.
\end{proof}

\begin{stwierdzenie}[Odwrotność własności średniej]
\label{prop:pde3-converse}
Niech $n\geq2$, $\Omega\subset\R^n$ będzie otwarta,
a $u\in C^2(\Omega)$. Jeśli dla każdego $x\in\Omega$
i wszystkich dostatecznie małych $r>0$ spełniona jest
pierwsza równość w~\eqref{eq:pde3-both-means}, to
$\Delta u=0$ w $\Omega$. Ten sam wniosek zachodzi przy
założeniu własności średniej po kulach.
\end{stwierdzenie}
\begin{proof}
Ustalmy $x\in\Omega$ i wybierzmy $r_0>0$ takie,
że $\overline{B_{r_0}(x)}\subset\Omega$.
Dla $|h|<r_0$ wzór Taylora z resztą całkową daje
\[
 \begin{aligned}
 u(x+h)&=u(x)+\nabla u(x)\cdot h
 +\tfrac12 h^{\mathsf T}D^2u(x)h+E(h),\\
 E(h)&=\int_0^1(1-t)h^{\mathsf T}
 \bigl(D^2u(x+th)-D^2u(x)\bigr)h\,\mathrm d\lambda_1(t).
 \end{aligned}
\]
Z ciągłości hesjanu wynika
$|E(h)|\leq |h|^2\eta(|h|)$, gdzie
$\eta(r)=\frac12\sup_{|z-x|\leq r}\|D^2u(z)-D^2u(x)\|$
dąży do zera przy $r\downarrow0$, a norma macierzy
jest normą operatorową indukowaną przez normę
euklidesową. W szczególności reszta dla $h=r\theta$
jest $o(r^2)$ jednostajnie względem $\theta\in S^{n-1}$.

Odbicie jednej współrzędnej zachowuje miarę sfery,
więc całki z $\theta_i$ oraz z $\theta_i\theta_j$
dla $i\ne j$ znikają. Permutacje współrzędnych dają
równość wszystkich całek z $\theta_i^2$.
Ponieważ $\sum_{i=1}^n\theta_i^2=1$ na sferze,
każda z tych całek wynosi $\omega_n/n$. Stąd
\[
 M_u(x,r)=u(x)+\frac{r^2}{2n}\Delta u(x)+e(r),
 \qquad |e(r)|\leq r^2\eta(r).
\]
Jeśli $M_u(x,r)=u(x)$ dla wszystkich dostatecznie małych
$r$, po odjęciu $u(x)$ i podzieleniu przez $r^2$
przechodzimy do granicy, otrzymując $\Delta u(x)=0$.

Dla średniej po kuli całkujemy ostatni wzór
w zmiennej promieniowej. Otrzymujemy
\[
 \begin{aligned}
 \frac{1}{\lambda_n(B_r(x))}
 \int_{B_r(x)}u(y)\,\mathrm d\lambda_n(y)
 &=\frac{n}{r^n}\int_0^r M_u(x,s)s^{n-1}\,
       \mathrm d\lambda_1(s)\\
 &=u(x)+\frac{r^2}{2(n+2)}\Delta u(x)+E_r .
 \end{aligned}
\]
Ponieważ $\eta$ jest niemalejąca,
$|E_r|\leq \frac{n}{n+2}r^2\eta(r)=o(r^2)$.
Własność średniej po kulach również daje więc
$\Delta u(x)=0$. Punkt $x$ był dowolny.
\end{proof}

\begin{uwaga}[Dlaczego średnia po kulach nie jest przypadkowa?]
Wartość funkcji harmonicznej w środku jest dokładnie
równowagą jej wartości wokół środka. Dla funkcji
subharmonicznej o klasie $C^2$ rachunek z dowodu pokazuje,
że $M_u(x,r)$ jest niemalejąca jako funkcja $r$, więc
$u(x)\leq M_u(x,r)$. Dla funkcji superharmonicznej
nierówność ma kierunek przeciwny.
\end{uwaga}

\begin{twierdzenie}[Własność wartości średniej bez założenia $C^2$]
\label{thm:pde-v31-continuous-mean}
Niech $n\ge2$, $\Omega\subset\R^n$ będzie otwarta,
a $u\in C(\Omega)$. Załóżmy, że dla każdego
$x\in\Omega$ i każdego $r>0$ z
$\overline{B_r(x)}\subset\Omega$ zachodzi
\[
 u(x)=\frac1{\omega_nr^{n-1}}
 \int_{\partial B_r(x)}u(y)\,\mathrm d\sigma(y).
\]
Wówczas $u\in C^\infty(\Omega)$ i $\Delta u=0$ w $\Omega$.
\end{twierdzenie}
\begin{proof}
Ustalmy $x_0\in\Omega$ i wybierzmy $R>0$ takie,
że $\overline{B_{2R}(x_0)}\subset\Omega$.
Wybierzmy radialną, nieujemną funkcję
$\rho\in C_c^\infty(B_1(0))$ o całce $1$,
na przykład odpowiednio unormowaną funkcję postaci
$\rho(z)=c\,\exp(-1/(1-4|z|^2))$ dla $|z|<1/2$
i $\rho(z)=0$ w pozostałych punktach.
Dla $0<\eps<R/2$ połóżmy
$\rho_\eps(z)=\eps^{-n}\rho(z/\eps)$.
Zmiana zmiennych $z=\eps w$ daje
$\int_{\R^n}\rho_\eps\,\mathrm d\lambda_n=1$.
Dla $x\in B_{3R/2}(x_0)$ określmy
\[
 u_\eps(x)=\int_{B_\eps(0)}\rho_\eps(z)u(x-z)\,
 \mathrm d\lambda_n(z).
\]
Argument $x-z$ należy do $B_{2R}(x_0)$,
więc całka jest dobrze określona. Zmieniając zmienne
$y=x-z$ i korzystając z nośnika $\rho_\eps$, mamy
\[
 u_\eps(x)=\int_{B_{2R}(x_0)}\rho_\eps(x-y)u(y)\,
 \mathrm d\lambda_n(y).
\]
Funkcja $u$ jest ograniczona na $\overline{B_{2R}(x_0)}$,
a każda pochodna $\rho_\eps$ jest ograniczona na $\R^n$.
Dla każdej pochodnej względem $x$ istnieje zatem
stała całkowalna na $B_{2R}(x_0)$, dominująca moduł
odpowiedniej funkcji podcałkowej.
Wielokrotne różniczkowanie pod całką daje
$u_\eps\in C^\infty(B_{3R/2}(x_0))$.

Radialność pozwala napisać $\rho_\eps(z)=k_\eps(|z|)$.
Dla $x\in B_{3R/2}(x_0)$ i $0<s<\eps$ mamy
$\overline{B_s(x)}\subset B_{2R}(x_0)$.
Stosując współrzędne biegunowe oraz założoną własność
średniej na sferze o środku $x$, otrzymujemy
\[
 \begin{aligned}
 u_\eps(x)
 &=\int_0^\eps k_\eps(s)s^{n-1}
       \int_{S^{n-1}}u(x-s\theta)\,\mathrm d\sigma(\theta)
       \,\mathrm d\lambda_1(s)\\
 &=u(x)\omega_n\int_0^\eps k_\eps(s)s^{n-1}
       \,\mathrm d\lambda_1(s)
 =u(x).
 \end{aligned}
\]
Zmiana $\theta\mapsto-\theta$ zachowuje miarę sfery,
a ostatnia równość wynika z normalizacji $\rho_\eps$.
Zatem $u$ jest gładka na $B_{3R/2}(x_0)$, bez potrzeby
przechodzenia z $\eps$ do zera. Stwierdzenie
\ref{prop:pde3-converse} daje $\Delta u=0$ w tym zbiorze.
Dowolność $x_0$ kończy dowód.

Warto zauważyć także inną drogę: dla dowolnego (niekoniecznie
radialnego) nieujemnego jądra wygładzającego o całce $1$
i nośniku w $B_\eps(0)$ powyższy splot jest gładki.
Dla $x\in B_R(x_0)$ oraz $0<r<R/2$ twierdzenie Fubiniego
można zastosować do całki po $\partial B_r(x)\times B_\eps(0)$,
ponieważ funkcja podcałkowa jest ograniczona, a obie miary
są skończone. Własność średniej $u$ dla kul o środkach
$x-z$ daje własność średniej $u_\eps$; stąd
$\Delta u_\eps=0$ na $B_R(x_0)$. Jednostajna ciągłość
$u$ na $\overline{B_{2R}(x_0)}$ daje
$u_\eps\to u$ jednostajnie na $\overline{B_R(x_0)}$.
Po udowodnieniu twierdzenia o jądrze Poissona
(\ref{thm:pde3-poisson}) można zastosować jego wzór
do $u_\eps$ i przejść do granicy, korzystając z dodatniości
i normalizacji jądra. Otrzymuje się w ten sposób
alternatywny dowód harmoniczności i gładkości $u$.
\end{proof}

\begin{uwaga}
Kluczowe jest, że własność średniej zakładamy dla
\emph{wszystkich} kul o domknięciu zawartym w $\Omega$.
Nie należy bez dodatkowego argumentu zastępować tego
warunku wymaganiem dotyczącym jednego promienia
zależnego od środka kuli.
\end{uwaga}

\section{Mocna zasada maksimum, jednoznaczność i twierdzenie Liouville'a}
\label{sec:pde3-max}

Słabą zasadę maksimum udowodniliśmy już w twierdzeniu
\ref{thm:laplace-max}. Własność wartości średniej daje
twierdzenie silniejsze, które nie wymaga ograniczoności obszaru.

\begin{twierdzenie}[Mocna zasada maksimum dla funkcji harmonicznych]
\label{thm:pde3-strong-max}
Niech $\Omega\subset\R^n$ będzie obszarem, $n\geq2$,
a $u\in C^2(\Omega)$ funkcją harmoniczną. Jeśli
$u$ osiąga maksimum lokalne albo minimum lokalne w pewnym
punkcie $x_0\in\Omega$, to jest stała na $\Omega$.
\end{twierdzenie}
\begin{proof}
Niech $u$ osiąga maksimum lokalne w $x_0$ i oznaczmy
$M=u(x_0)$. Wybierzmy $r_0>0$ tak, aby
$\overline{B_{r_0}(x_0)}\subset\Omega$ oraz $u\leq M$
na $B_{r_0}(x_0)$. Dla każdego $0<r<r_0$ własność
wartości średniej daje
\[
 \int_{\partial B_r(x_0)}(M-u(y))\,\mathrm d\sigma(y)=0.
\]
Funkcja podcałkowa jest ciągła i nieujemna.
Gdyby była dodatnia w pewnym punkcie sfery,
byłaby ograniczona z dołu przez dodatnią stałą
na pewnym niepustym względnie otwartym fragmencie
sfery. Każdy taki fragment ma dodatnią miarę
powierzchniową, co przeczyłoby zerowaniu się całki.
Stąd $u=M$ na każdej sferze $\partial B_r(x_0)$
dla $0<r<r_0$, a zatem na całej kuli $B_{r_0}(x_0)$.

Jeśli $M$ jest maksimum globalnym, można zakończyć
dowód bez analityczności: zbiór $\{x\in\Omega:u(x)=M\}$
jest niepusty i względnie domknięty z ciągłości $u$.
Powyższy argument zastosowany w dowolnym jego punkcie
pokazuje, że jest również otwarty. Spójność $\Omega$
wymusza jego równość z $\Omega$.

Dla maksimum jedynie lokalnego nie możemy zakładać
nierówności $u\leq M$ na całym obszarze.
Korzystamy wówczas z twierdzenia o analityczności
funkcji harmonicznych (twierdzenie~\ref{thm:pde3-analytic}).
Niech $w=u-M$ i niech $\mathbb N_0=\{0,1,2,\ldots\}$.
Rozważmy zbiór
\[
 A=\{x\in\Omega:D^\alpha w(x)=0
       \text{ dla każdego }\alpha\in\mathbb N_0^n\},
\]
gdzie $D^\alpha=\partial_{x_1}^{\alpha_1}\cdots
\partial_{x_n}^{\alpha_n}$ i $D^0w=w$.
Zbiór $A$ zawiera $B_{r_0}(x_0)$ i jest względnie
domknięty jako część wspólna zbiorów zerowych
ciągłych pochodnych $D^\alpha w$.
Jeżeli $x\in A$, to rzeczywista analityczność $w$
oznacza, że na pewnej kuli o środku $x$ funkcja $w$
jest sumą swojego szeregu Taylora. Wszystkie jego
współczynniki $D^\alpha w(x)/\alpha!$ są zerowe
(tutaj $\alpha!=\prod_{j=1}^n\alpha_j!$).
Stąd $w$ znika na tej kuli, a wraz z nią wszystkie
jej pochodne. Kula ta zawiera się więc w $A$.
Wynika stąd, że $A$ jest również otwarty.
Ze spójności $\Omega$ otrzymujemy $A=\Omega$,
a w szczególności $w=0$ na $\Omega$.

Nie ma tu błędnego koła: dowód analityczności korzysta
z twierdzenia o jądrze Poissona, którego jednoznaczność
wynika ze słabej zasady maksimum udowodnionej w rozdziale
pierwszym, a konstrukcja i normalizacja jądra korzystają
z funkcji Greena i wzoru reprezentacyjnego Greena.
Żaden z tych argumentów nie wymaga obecnego twierdzenia.
Przypadek minimum sprowadza się do maksimum funkcji $-u$.
\end{proof}

\begin{stwierdzenie}[Oszacowanie pochodnych w środku kuli]
\label{prop:pde3-gradient}
Niech $n\geq2$, $R>0$, $x_0\in\R^n$,
a $u\in C^2(B_R(x_0))\cap C(\overline{B_R(x_0)})$
będzie harmoniczna. Wówczas dla każdego
$j=1,\dots,n$ zachodzi
\begin{equation}\label{eq:pde3-derivative-bound}
 |\partial_{x_j}u(x_0)|\leq \frac nR
 \sup_{y\in B_R(x_0)}|u(y)|.
\end{equation}
\end{stwierdzenie}
\begin{proof}
Ustalmy $0<\rho<R$ i połóżmy $\delta=(R-\rho)/2$.
Jeśli $x\in B_\delta(x_0)$, to dla każdego
$z\in\overline{B_\rho(0)}$ mamy
$|x+z-x_0|\leq|x-x_0|+|z|<\delta+\rho<R$.
Możemy więc zastosować własność wartości średniej
po kuli o ustalonym promieniu $\rho$ i ruchomym środku $x$:
\[
 u(x)=\frac{1}{\lambda_n(B_\rho(0))}
 \int_{B_\rho(0)}u(x+z)\,\mathrm d\lambda_n(z).
\]
Funkcja podcałkowa oraz jej pochodne względem $x$
są ciągłe i ograniczone dla
$(x,z)\in\overline{B_\delta(x_0)}\times
\overline{B_\rho(0)}$, ponieważ wszystkie argumenty $x+z$
leżą w zwartej kuli o promieniu $\delta+\rho<R$.
Różniczkując względem $x_j$ pod znakiem całki
w punkcie $x_0$ i podstawiając $y=x_0+z$, otrzymujemy
\[
 \partial_{x_j}u(x_0)
 =\frac{1}{\lambda_n(B_\rho(0))}
 \int_{B_\rho(x_0)}\partial_{x_j}u(y)\,
 \mathrm d\lambda_n(y).
\]
Nie potrzebujemy przy tym zakładać harmoniczności
pochodnych $u$, która wymagałaby dodatkowej regularności.
Dla $e_j$ oznaczającego $j$-ty wektor standardowej
bazy $\R^n$ pole $y\mapsto u(y)e_j$ jest klasy $C^1$
na otoczeniu domkniętej kuli, a jego dywergencją jest
$\partial_{x_j}u$. Twierdzenie Gaussa daje zatem
\[
 \partial_{x_j}u(x_0)=\frac{1}{\lambda_n(B_\rho(0))}
 \int_{\partial B_\rho(x_0)}u(y)\nu_j(y)\,
 \mathrm d\sigma(y),
\]
gdzie $\nu_j(y)$ jest $j$-tą współrzędną jednostkowej
normalnej zewnętrznej $\nu(y)=(y-x_0)/\rho$.
Oznaczmy $K=\sup_{y\in B_R(x_0)}|u(y)|<\infty$,
co wynika z ciągłości $u$ na domkniętej kuli.
Korzystając z $|\nu_j|\leq1$, otrzymujemy
\[
 |\partial_{x_j}u(x_0)|
 \leq K\frac{\omega_n\rho^{n-1}}{\omega_n\rho^n/n}
 =\frac{nK}{\rho}.
\]
Lewa strona i $K$ nie zależą od $\rho$. Przechodząc
z $\rho$ do $R$ od dołu, dostajemy
\eqref{eq:pde3-derivative-bound}.
\end{proof}

\begin{wniosek}[Twierdzenie Liouville'a]
\label{cor:pde3-liouville}
Niech $n\geq2$ i $u\in C^2(\R^n)$ będzie harmoniczna
oraz ograniczona. Wtedy $u$ jest stała.
\end{wniosek}
\begin{proof}
Niech $K=\sup_{x\in\R^n}|u(x)|<\infty$.
Dla każdego $x_0\in\R^n$, $R>0$ oraz $j=1,\ldots,n$
możemy zastosować \eqref{eq:pde3-derivative-bound},
gdyż $u$ jest harmoniczna na $B_R(x_0)$
i ciągła na jej domknięciu. Otrzymujemy
$|\partial_{x_j}u(x_0)|\leq nK/R$.
Przechodząc z $R$ do nieskończoności przy ustalonym $x_0$,
dostajemy $\partial_{x_j}u(x_0)=0$ dla każdego $j$.
Ponieważ $x_0$ był dowolny, mamy $\nabla u=0$ na $\R^n$.
Dla dowolnych $x,y\in\R^n$ reguła łańcuchowa
i twierdzenie podstawowe rachunku całkowego dają
\[
 u(y)-u(x)=\int_0^1
 \nabla u(x+t(y-x))\cdot(y-x)\,\mathrm d\lambda_1(t)=0.
\]
Zatem $u$ jest stała.
\end{proof}

\begin{kontrprzyklad}[Dlaczego ograniczoność jest ważna?]
Funkcja $u(x)=x_1$ jest harmoniczna na $\R^n$,
lecz nie jest stała. Twierdzenie Liouville'a nie
obowiązuje bez ograniczoności (ani analogicznej
kontroli wzrostu).
\end{kontrprzyklad}

\section{Rozwiązanie fundamentalne laplasjanu}
\label{sec:pde3-fundamental}

Rozwiązaniem fundamentalnym operatora nazywa się funkcję lub dystrybucję,
która reaguje na punktowe źródło. Wprowadzimy ją najpierw
bez formalnej teorii dystrybucji: zapis
$-\Delta\Phi=\delta_0$ będzie oznaczał precyzyjną
tożsamość całkową, gdzie $\delta_0$ jest funkcjonałem
$\varphi\mapsto\varphi(0)$ na funkcjach testowych.
Tutaj $C_c^\infty(\R^n)$ oznacza przestrzeń funkcji
nieskończenie różniczkowalnych o zwartym nośniku, a
$L^1_{\mathrm{loc}}(\R^n)$ klasę funkcji całkowalnych
na każdym zwartym podzbiorze $\R^n$.

\begin{twierdzenie}[Rozwiązanie fundamentalne $-\Delta$]
\label{thm:pde3-phi}
Niech $n\geq2$ i $\omega_n=\sigma(S^{n-1})$. Dla
$x\in\R^n\setminus\{0\}$ połóżmy
\begin{equation}\label{eq:pde3-fundamental}
 \Phi_n(x)=
 \begin{cases}
 \displaystyle\frac{1}{(n-2)\omega_n}|x|^{2-n},&n\geq3,\\[4pt]
 \displaystyle-\frac1{2\pi}\log|x|,&n=2.
 \end{cases}
\end{equation}
Wartość $\Phi_n(0)$ ustalamy dowolnie, na przykład
$\Phi_n(0)=0$; nie wpływa ona na poniższe całki.
Wówczas $\Phi_n\in L^1_{\mathrm{loc}}(\R^n)$,
$\Delta\Phi_n=0$ na $\R^n\setminus\{0\}$ oraz
\begin{equation}\label{eq:pde3-distdelta}
 \int_{\R^n}\Phi_n(x)(-\Delta\varphi(x))\,
 \mathrm d\lambda_n(x)=\varphi(0)
 \quad\text{dla każdej }\varphi\in C_c^\infty(\R^n).
\end{equation}
\end{twierdzenie}
\begin{proof}
Niech $h\in C^2((0,\infty))$ i $w(x)=h(|x|)$
dla $x\ne0$. Oznaczmy $r=|x|$. Ponieważ
$\partial_{x_j}r=x_j/r$, reguła łańcuchowa daje
\[
 \partial_{x_j}w=h'(r)\frac{x_j}{r},\qquad
 \partial_{x_j}^2w
 =h''(r)\frac{x_j^2}{r^2}
 +h'(r)\left(\frac1r-\frac{x_j^2}{r^3}\right).
\]
Sumując względem $j$ i korzystając z
$\sum_{j=1}^n x_j^2=r^2$, otrzymujemy
$\Delta w=h''(r)+(n-1)h'(r)/r$.
Równanie $\Delta w=0$ jest więc równoważne
$(r^{n-1}h'(r))'=0$, skąd $h'(r)=Cr^{1-n}$.
Całkowanie daje $h(r)=A+Cr^{2-n}/(2-n)$ dla $n\geq3$
oraz $h(r)=A+C\log r$ dla $n=2$.

Aby otrzymać jednostkowe punktowe źródło dla operatora
$-\Delta$, wybieramy $C=-1/\omega_n$; wtedy strumień
$-\nabla w$ przez każdą sferę o środku w zerze wynosi
$-h'(r)\omega_n r^{n-1}=1$.
Przy $A=0$ i $\omega_2=2\pi$ otrzymujemy dokładnie
\eqref{eq:pde3-fundamental}. Oznaczmy przez $h_n(r)$
profil radialny tej funkcji, tak że
$\Phi_n(x)=h_n(|x|)$ dla $x\ne0$. Mamy
$h_n'(r)=-1/(\omega_n r^{n-1})$ i
$\nabla\Phi_n(x)=-x/(\omega_n |x|^n)$.
Wykażemy poniżej, że wybrany strumień daje również
żądaną tożsamość całkową.

Funkcja $\Phi_n$ jest gładka poza zerem.
Dla $n\geq3$ i $a>0$ współrzędne biegunowe dają
\[
 \int_{B_a(0)}|\Phi_n(x)|\,\mathrm d\lambda_n(x)
 =\frac1{n-2}\int_0^a r\,\mathrm d\lambda_1(r)
 =\frac{a^2}{2(n-2)}.
\]
Dla $n=2$ oraz $0<a\leq1$ analogicznie
\[
 \int_{B_a(0)}|\Phi_2(x)|\,\mathrm d\lambda_2(x)
 =\int_0^a(-\log r)r\,\mathrm d\lambda_1(r)
 =\frac{a^2}{4}-\frac{a^2}{2}\log a.
\]
W ostatnim rachunku używamy
$\lim_{r\downarrow0}r^2\log r=0$.
Obie całki są skończone i dążą do zera wraz z $a$.
W połączeniu z gładkością poza zerem daje to
$\Phi_n\in L^1_{\mathrm{loc}}(\R^n)$.

Niech $\varphi\in C_c^\infty(\R^n)$ i wybierzmy $R>0$
tak, aby jej nośnik był zawarty w $B_R(0)$.
Dla $0<\eps<R$ rozważmy pierścień
$A_{\eps,R}=B_R(0)\setminus\overline{B_\eps(0)}$.
Funkcje $\Phi_n$ i $\varphi$ są gładkie na otoczeniu
jego domknięcia. Stosujemy twierdzenie Gaussa do pola
$\varphi\nabla\Phi_n-\Phi_n\nabla\varphi$,
którego dywergencja jest równa
$\varphi\Delta\Phi_n-\Phi_n\Delta\varphi
=-\Phi_n\Delta\varphi$ na pierścieniu.
Jest to druga tożsamość Greena, wyprowadzona tu
bez odwoływania się do dalszej części rozdziału.
Na zewnętrznym brzegu $\varphi$ i jej pochodne
znikają, gdyż $\varphi=0$ w jego otoczeniu.
Na wewnętrznym brzegu jednostkowa normalna zewnętrzna
dla pierścienia wynosi $\nu(y)=-y/\eps$,
więc $\partial_\nu\Phi_n(y)=-h_n'(\eps)
=1/(\omega_n\eps^{n-1})$.
Otrzymujemy
\[
 \int_{A_{\eps,R}}\Phi_n(y)(-\Delta\varphi(y))\,
 \mathrm d\lambda_n(y)
 =\int_{\partial B_\eps(0)}
 \bigl(\varphi(y)\,\partial_\nu\Phi_n(y)
       -\Phi_n(y)\,\partial_\nu\varphi(y)\bigr)\,\mathrm d\sigma(y).
\]
Pierwszy składnik całki brzegowej jest średnią
sferyczną funkcji $\varphi$:
\[
 \left|\frac1{\omega_n\eps^{n-1}}
 \int_{\partial B_\eps(0)}\varphi(y)\,\mathrm d\sigma(y)
 -\varphi(0)\right|
 \leq\sup_{|y|=\eps}|\varphi(y)-\varphi(0)|
 \longrightarrow0.
\]
Dla drugiego składnika, z $K=\sup_{\R^n}|\nabla\varphi|$,
mamy oszacowanie
\[
 \left|\int_{\partial B_\eps(0)}
 \Phi_n(y)\partial_\nu\varphi(y)\,\mathrm d\sigma(y)\right|
 \leq K\omega_n\eps^{n-1}|h_n(\eps)|.
\]
Prawa strona wynosi $K\eps/(n-2)$ dla $n\geq3$,
a dla $n=2$ jest równa $K\eps|\log\eps|$.
W obu przypadkach dąży do zera.
Wreszcie całka po $\R^n$ różni się od całki po
$A_{\eps,R}$ tylko całką po $B_\eps(0)$ (z dokładnością
do zbiorów miary zero), ponieważ $\Delta\varphi=0$
poza $B_R(0)$. Moduł tej różnicy jest nie większy niż
$\sup_{\R^n}|\Delta\varphi|
\int_{B_\eps(0)}|\Phi_n(y)|\,\mathrm d\lambda_n(y)$,
co dąży do zera na mocy wcześniejszych obliczeń.
Przechodząc do granicy otrzymujemy
\eqref{eq:pde3-distdelta}.
\end{proof}

\begin{uwaga}[Wymiar jeden]
Dla $n=1$ odpowiednikiem rozwiązania fundamentalnego
jest $\Phi_1(x)=-|x|/2$. Istotnie, dla dowolnej
$\varphi\in C_c^\infty(\R)$ całkowanie przez części
na obu półprostych daje
\[
 \begin{aligned}
 \int_{-\infty}^0(-x)\varphi''(x)\,\mathrm d\lambda_1(x)
 &=\varphi(0),\\
 \int_0^\infty x\varphi''(x)\,\mathrm d\lambda_1(x)
 &=\varphi(0).
 \end{aligned}
\]
Wyrazy brzegowe w nieskończoności znikają dzięki
zwartemu nośnikowi, a iloczyn $x\varphi'(x)$ znika
w zerze. Stąd $\int_\R \Phi_1(-\varphi'')\,
\mathrm d\lambda_1=\varphi(0)$, czyli
$-\Phi_1''=\delta_0$ w sensie dystrybucji.
Przypadek jednowymiarowy jest w istocie teorią ODE,
dlatego w dalszych konstrukcjach potencjału przyjmujemy $n\geq2$.
\end{uwaga}

\begin{twierdzenie}[Potencjał Newtona i równanie Poissona]
\label{thm:pde3-newton}
Niech $n\geq2$ i $f\in C_c^\infty(\R^n)$. Określmy
\begin{equation}\label{eq:pde3-newton}
 (\Phi_n*f)(x)=\int_{\R^n}\Phi_n(x-y)f(y)\,
 \mathrm d\lambda_n(y).
\end{equation}
Wówczas $u=\Phi_n*f\in C^\infty(\R^n)$ i
$-\Delta u=f$ na $\R^n$.
\end{twierdzenie}
\begin{proof}
Ustalmy $x_0\in\R^n$ i zwartą kulę
$K=\overline{B_1(x_0)}$. Niech $L=\supp f$;
dla $f=0$ teza jest oczywista, więc zakładamy $L\ne\emptyset$.
Zbiór $S=K-L=\{x-y:x\in K,\ y\in L\}$ jest zwarty
jako ciągły obraz zbioru zwartego $K\times L$.
Zmiana zmiennych $z=x-y$ ma moduł wyznacznika równy $1$.
Ponieważ $f$ jest ograniczona, a $\Phi_n$ całkowalna
na $S$, dla każdego $x\in K$ całka definiująca $u$ jest
bezwzględnie zbieżna i
\[
 u(x)=\int_{\R^n}\Phi_n(z)f(x-z)\,\mathrm d\lambda_n(z).
\]
Dla każdego wielowskaźnika $\alpha\in\mathbb N_0^n$
pochodna $D^\alpha f$ ma nośnik zawarty w $L$,
a $M_\alpha=\sup_{\R^n}|D^\alpha f|<\infty$.
Wobec tego dla $x\in K$ mamy
\[
 |\Phi_n(z)D^\alpha f(x-z)|
 \leq M_\alpha|\Phi_n(z)|\chi_S(z).
\]
Prawa strona jest całkowalna względem $\lambda_n$.
Na wnętrzu $K$ można więc wielokrotnie różniczkować
pod całką: na każdym kroku ilorazy różnicowe
są zdominowane przez powyższe oszacowanie dla
odpowiedniej pochodnej następnego rzędu, z twierdzenia
o wartości średniej, a ciągłość otrzymanych całek
wynika z twierdzenia o zbieżności zdominowanej.
Dostajemy
\[
 D^\alpha u(x)=\int_{\R^n}
 \Phi_n(z)D^\alpha f(x-z)\,\mathrm d\lambda_n(z)
 \qquad(x\in B_1(x_0)).
\]
Ponieważ $x_0$ był dowolny, $u\in C^\infty(\R^n)$.

Dla ustalonego $x\in\R^n$ funkcja
$\varphi_x(z)=f(x-z)$ należy do $C_c^\infty(\R^n)$.
Reguła łańcuchowa daje
$\partial_{z_j}\varphi_x(z)=-\partial_{x_j}f(x-z)$,
a przy drugim różniczkowaniu znak minus znika.
Zatem $\Delta_z\varphi_x(z)=(\Delta f)(x-z)$.
Stosując \eqref{eq:pde3-distdelta} do $\varphi_x$
i korzystając z uzasadnionego różniczkowania splotu,
otrzymujemy
\[
 -\Delta u(x)=\int_{\R^n}
 \Phi_n(z)(-\Delta_z\varphi_x(z))\,\mathrm d\lambda_n(z)
 =\varphi_x(0)=f(x).
\]
Równanie zachodzi punktowo na całej przestrzeni.
\end{proof}

\begin{uwaga}[Rozwiązanie szczególne a niejednoznaczność]
Jeżeli $u,v\in C^2(\R^n)$ spełniają
$-\Delta u=-\Delta v=f$, to ich różnica jest harmoniczna.
Odwrotnie, dodanie dowolnej funkcji harmonicznej
na $\R^n$ do rozwiązania zachowuje równanie.
Zatem dla $f\in C_c^\infty(\R^n)$ wszystkie rozwiązania
klasyczne mają postać $\Phi_n*f+h$, gdzie
$h\in C^2(\R^n)$ i $\Delta h=0$.
Przykładowo $\Phi_n*f+x_1$ jest innym rozwiązaniem.
Do jednoznaczności potrzebne są dodatkowe warunki
eliminujące swobodę dodawania funkcji harmonicznych.
\end{uwaga}

\section{Funkcje Greena i tożsamości Greena}
\label{sec:pde3-green}

\begin{stwierdzenie}[Dwie tożsamości Greena]
\label{prop:pde3-green-identities}
Niech $\Omega\subset\R^n$ będzie obszarem ograniczonym
o brzegu klasy $C^1$ i niech $u,v\in C^2(\overline\Omega)$
będą funkcjami rzeczywistymi. Przez $C^2(\overline\Omega)$
rozumiemy tutaj funkcje klasy $C^2$ w $\Omega$, których
wartości i pochodne do rzędu drugiego przedłużają się
ciągle na $\overline\Omega$. Niech $\nu$ oznacza
jednostkową normalną zewnętrzną do $\partial\Omega$,
a $\partial_\nu v=\nabla v\cdot\nu$. Wtedy
\begin{align}
 \int_\Omega u(x)\Delta v(x)\,\mathrm d\lambda_n(x)
 &=-\int_\Omega\nabla u(x)\cdot\nabla v(x)\,\mathrm d\lambda_n(x)
 +\int_{\partial\Omega}u(x)\,\partial_\nu v(x)\,\mathrm d\sigma(x),
 \label{eq:pde3-green1}\\
 \int_\Omega(u(x)\Delta v(x)-v(x)\Delta u(x))\,\mathrm d\lambda_n(x)
 &=\int_{\partial\Omega}(u(x)\partial_\nu v(x)-v(x)\partial_\nu u(x))
 \,\mathrm d\sigma(x).\label{eq:pde3-green2}
\end{align}
\end{stwierdzenie}
\begin{proof}
Pole wektorowe $u\nabla v$ jest klasy $C^1$
aż do brzegu, a jego dywergencja wynosi
$\diver(u\nabla v)=\nabla u\cdot\nabla v+u\Delta v$.
Twierdzenie Gaussa daje
\[
 \int_\Omega(\nabla u\cdot\nabla v+u\Delta v)\,
 \mathrm d\lambda_n
 =\int_{\partial\Omega}u\,\partial_\nu v\,\mathrm d\sigma.
\]
Wszystkie całki są skończone, ponieważ funkcje
podcałkowe są ciągłe na odpowiednich zbiorach
zwartych, a ograniczony obszar o brzegu klasy $C^1$
ma skończoną objętość i skończoną miarę brzegu.
Przeniesienie składnika z gradientami daje
\eqref{eq:pde3-green1}. Zapisując tę samą tożsamość
dla pary $(v,u)$ i odejmując ją od pierwszej,
korzystamy z równości $\nabla u\cdot\nabla v
=\nabla v\cdot\nabla u$. Otrzymujemy
\eqref{eq:pde3-green2}.
\end{proof}

\begin{definicja}[Funkcja Greena dla $-\Delta$]
\label{def:pde3-green}
Niech $\Omega\subset\R^n$ będzie obszarem ograniczonym,
$n\geq2$. Dla ustalonego $x\in\Omega$ funkcją Greena
(z warunkiem Dirichleta) nazywamy funkcję
$G_\Omega(x,\cdot)$ taką, że:
\begin{enumerate}
\item jest lokalnie całkowalna w $\Omega$ i klasy
$C^2$ na $\Omega\setminus\{x\}$;
\item $-\Delta_yG_\Omega(x,y)=\delta_x$ w sensie
dystrybucji, to znaczy
\[
 \int_\Omega G_\Omega(x,y)(-\Delta\varphi(y))\,
 \mathrm d\lambda_n(y)=\varphi(x)
 \quad(\varphi\in C_c^\infty(\Omega));
\]
\item funkcja $G_\Omega(x,\cdot)$ ma ciągłe przedłużenie
na $\overline\Omega\setminus\{x\}$, którego wartości
na $\partial\Omega$ są równe zeru.
\end{enumerate}
Wartość w punkcie $y=x$ nie ma znaczenia dla tożsamości
całkowej. Drugi warunek wraz z pierwszym wymusza
$\Delta_yG_\Omega(x,y)=0$ dla $y\ne x$: dla funkcji
testowych o nośniku z dala od $x$ dwukrotne całkowanie
przez części daje zerową całkę z ciągłej funkcji
$\Delta_yG_\Omega$ pomnożonej przez dowolną taką funkcję
testową. Gdyby ta ciągła funkcja była niezerowa w pewnym
punkcie, test nieujemny o nośniku w otoczeniu, na którym
ma ona stały niezerowy znak, dawałby sprzeczność.
W zastosowaniach do całek brzegowych wymagamy również
odpowiedniej gładkości funkcji Greena względem $y$ aż do brzegu.
\end{definicja}

\begin{twierdzenie}[Wzór reprezentacyjny Greena]
\label{thm:pde3-representation}
Niech $n\geq2$, $\Omega\subset\R^n$ będzie obszarem ograniczonym o brzegu klasy $C^2$, $f\in C(\overline\Omega)$
i $u\in C^2(\overline\Omega)$ spełnia $-\Delta u=f$
w $\Omega$. Załóżmy, że dla każdego $x\in\Omega$
istnieje funkcja Greena $G_\Omega(x,y)=\Phi_n(x-y)-h_x(y)$,
gdzie $h_x$ jest harmoniczna i klasy $C^2$ na pewnym
otoczeniu $\overline\Omega$ (w zmiennej $y$).
Wówczas dla $x\in\Omega$
\begin{equation}\label{eq:pde3-green-representation}
 u(x)=\int_\Omega G_\Omega(x,y)f(y)\,\mathrm d\lambda_n(y)
 -\int_{\partial\Omega}u(y)\partial_{\nu_y}G_\Omega(x,y)
 \,\mathrm d\sigma(y).
\end{equation}
\end{twierdzenie}
\begin{proof}
Ustalmy $x\in\Omega$ i wybierzmy
$0<\eps<\operatorname{dist}(x,\partial\Omega)$.
Dla skrócenia zapisu połóżmy $G(y)=G_\Omega(x,y)$.
Funkcje $u$ i $G$ są klasy $C^2$ aż do brzegu
zbioru $\Omega_\eps=\Omega\setminus\overline{B_\eps(x)}$;
jego brzeg składa się z $\partial\Omega$ i $\partial B_\eps(x)$.
Druga tożsamość Greena dla $u$ i $G_\Omega(x,\cdot)$
na $\Omega_\eps=\Omega\setminus\overline{B_\eps(x)}$
daje, ponieważ $\Delta G_\Omega=0$ poza $x$,
\[
 \int_{\Omega_\eps}G_\Omega(x,y)f(y)\,\mathrm d\lambda_n(y)
 =\int_{\partial\Omega_\eps}
 \bigl(u(y)\partial_{\nu_y}G_\Omega(x,y)
 -G_\Omega(x,y)\partial_{\nu_y}u(y)\bigr)
 \,\mathrm d\sigma(y).
\]
Na $\partial\Omega$ normalna jest normalną zewnętrzną
do $\Omega$ i $G=0$, więc pozostaje całka
$\int_{\partial\Omega}u(y)\partial_{\nu_y}G(y)\,
\mathrm d\sigma(y)$.
Na $\partial B_\eps(x)$ normalna zewnętrzna
do $\Omega_\eps$ wynosi $\nu_y=-(y-x)/\eps$.
Z radialności $\Phi_n$ i rachunku w twierdzeniu
\ref{thm:pde3-phi} mamy
$\partial_{\nu_y}\Phi_n(x-y)=1/(\omega_n\eps^{n-1})$.
Wkład wewnętrznej sfery jest więc równy
\[
 \begin{aligned}
 I_\eps={}&\frac1{\omega_n\eps^{n-1}}
 \int_{\partial B_\eps(x)}u(y)\,\mathrm d\sigma(y)\\
 &-\int_{\partial B_\eps(x)}
 \Phi_n(x-y)\partial_{\nu_y}u(y)\,\mathrm d\sigma(y)\\
 &+\int_{\partial B_\eps(x)}
 (h_x(y)\partial_{\nu_y}u(y)
 -u(y)\partial_{\nu_y}h_x(y))\,\mathrm d\sigma(y).
 \end{aligned}
\]
Pierwszy składnik dąży do $u(x)$ z ciągłości $u$.
Wybierzmy $a>0$ z $\overline{B_a(x)}\subset\Omega$
i ograniczmy się do $\eps<a$.
Moduł drugiego składnika jest nie większy niż
$\sup_{\overline{B_a(x)}}|\nabla u|\,
\omega_n\eps^{n-1}|h_n(\eps)|$, gdzie $h_n$
oznacza profil radialny rozwiązania fundamentalnego.
Oszacowanie to dąży do zera, jak wykazano
w twierdzeniu~\ref{thm:pde3-phi}.
Moduł trzeciego składnika jest ograniczony przez
$C_x\omega_n\eps^{n-1}$, gdzie można przyjąć
$C_x=\sup_{\overline{B_a(x)}}(|h_x||\nabla u|
+|u||\nabla h_x|)<\infty$.
Ponieważ $n\geq2$, ten składnik również znika.
Zatem $I_\eps\to u(x)$.

Pozostaje przejście do granicy w całce objętościowej.
Funkcja $f$ jest ograniczona, a $G$ jest całkowalna
w $\Omega$: w otoczeniu $x$ wynika to z lokalnej
całkowalności $\Phi_n$ i ograniczoności $h_x$,
poza tym otoczeniem z ciągłości na zbiorze zwartym.
W szczególności moduł pominiętej całki po $B_\eps(x)$
jest nie większy niż
\[
 \sup_{\overline\Omega}|f|
 \left(\int_{B_\eps(0)}|\Phi_n(z)|\,\mathrm d\lambda_n(z)
 +\sup_{\overline{B_a(x)}}|h_x|\,\lambda_n(B_\eps(0))\right)
 \longrightarrow0.
\]
Ostatecznie otrzymujemy
\[
 \int_\Omega G_\Omega(x,y)f(y)\,\mathrm d\lambda_n(y)
 =\int_{\partial\Omega}u(y)\partial_{\nu_y}G_\Omega(x,y)
 \,\mathrm d\sigma(y)+u(x).
\]
Przeniesienie całki brzegowej daje
\eqref{eq:pde3-green-representation}.
\end{proof}

\section{Jawne jądro Poissona i zagadnienie Dirichleta na kuli}
\label{sec:pde3-poisson-kernel}

W tym podrozdziale $n\geq2$, $R>0$,
$B_R=B_R(0)$ i $\zeta\in\partial B_R$.
Dla $y\in B_R\setminus\{0\}$ wprowadzamy punkt odwrócony
$y^*=R^2y/|y|^2$, leżący poza kulą. Zachodzi tożsamość
\begin{equation}\label{eq:pde3-kelvin-identity}
 |z-y^*|=\frac R{|y|}|z-y| \qquad (|z|=R).
\end{equation}
Istotnie, $|y^*|=R^2/|y|>R$, a dla $|z|=R$
\[
 \begin{aligned}
 |z-y^*|^2
 &=R^2-2R^2\frac{z\cdot y}{|y|^2}
       +\frac{R^4}{|y|^2}\\
 &=\frac{R^2}{|y|^2}(|z|^2-2z\cdot y+|y|^2).
 \end{aligned}
\]
Pierwiastkując obie strony, otrzymujemy
\eqref{eq:pde3-kelvin-identity}.

\begin{stwierdzenie}[Funkcja Greena dla kuli]
\label{prop:pde3-ball-green}
Dla $n\geq3$ i $x,y\in B_R$, $x\ne y$, $y\ne0$, funkcja
\begin{equation}\label{eq:pde3-ball-green}
 G_R(x,y)=\Phi_n(x-y)-\Bigl(\frac R{|y|}\Bigr)^{n-2}
 \Phi_n(x-y^*)
\end{equation}
jest funkcją Greena dla kuli w zmiennej $x$.
Dla $n=2$ odpowiednia funkcja ma postać
\[
 G_R(x,y)=\frac1{2\pi}\log\frac{|y|\,|x-y^*|}{R\,|x-y|}.
\]
W obu przypadkach formuły przedłużają się dla $y=0$
oraz $x\in B_R\setminus\{0\}$ przez granicę (dla $n\geq3$:
$G_R(x,0)=\Phi_n(x)-1/((n-2)\omega_n R^{n-2})$;
dla $n=2$: $G_R(x,0)=(2\pi)^{-1}\log(R/|x|)$).
Ponadto $G_R(x,y)=G_R(y,x)$ dla wszystkich
$x,y\in B_R$ z $x\ne y$. W szczególności jest to
funkcja Greena również w zmiennej $y$ przy ustalonym $x$.
\end{stwierdzenie}
\begin{proof}
Dla $y\ne0$ i $x\ne y$ wprowadźmy pomocnicze oznaczenia
\[
 Q_R(x,y)=R^2-2x\cdot y+\frac{|x|^2|y|^2}{R^2},
 \qquad H_R(x,y)=h_n\bigl(\sqrt{Q_R(x,y)}\bigr),
\]
gdzie $h_n$ jest profilem radialnym $\Phi_n$.
Bezpośrednie rozwinięcie kwadratu odległości daje
$Q_R(x,y)=|y|^2|x-y^*|^2/R^2$.
Dla $n\geq3$ wynika stąd
$H_R(x,y)=(R/|y|)^{n-2}\Phi_n(x-y^*)$.
Dla $n=2$ mamy
$H_R(x,y)=-(2\pi)^{-1}\log(|y|\,|x-y^*|/R)$.
W obu przypadkach podany wzór przyjmuje zatem postać
$G_R(x,y)=\Phi_n(x-y)-H_R(x,y)$.

Z nierówności Cauchy'ego--Schwarza wynika
$Q_R(x,y)\geq(R-|x||y|/R)^2>0$,
gdy $|x||y|<R^2$. Dlatego $H_R$ jest dobrze określona
i klasy $C^\infty$ na tym otwartym zbiorze par $(x,y)$,
również dla $y=0$ i na przekątnej $x=y$ wewnątrz kuli.
W szczególności $H_R(x,0)=h_n(R)$, co daje obie
podane granice dla ustalonego $x\ne0$.
Jeżeli $0\leq r<R$, to dla $|x|\leq R$ i $|y|\leq r$
zachodzi $Q_R(x,y)\geq(R-r)^2>0$.
Zatem $H_R$ oraz wszystkie jej pochodne są ograniczone
na tym zwartym zbiorze. Własność ta będzie potrzebna
przy całkowaniu funkcji Greena względem bieguna.

Przy ustalonym $y\in B_R\setminus\{0\}$ funkcja
$x\mapsto H_R(x,y)$ jest harmoniczna na otoczeniu
$\overline{B_R}$: dla $n\geq3$ jest wielokrotnością
rozwiązania fundamentalnego z biegunem $y^*$ poza
domkniętą kulą, a dla $n=2$ różni się od
$\Phi_2(x-y^*)$ o stałą niezależną od $x$.
Dla $y=0$ jest stała, więc również harmoniczna.
Funkcja $\Phi_n(x-y)$ jest lokalnie całkowalna
i spełnia $-\Delta_x\Phi_n(x-y)=\delta_y$,
na mocy twierdzenia~\ref{thm:pde3-phi} i zmiany zmiennych.
Odejmowanie $H_R$ zachowuje tę tożsamość.

Dla $|x|=R$ mamy $Q_R(x,y)=|x-y|^2>0$,
więc $H_R(x,y)=\Phi_n(x-y)$ i $G_R(x,y)=0$.
Ciągłość aż do brzegu poza biegunem wynika z jawnych
wzorów. Wszystkie warunki definicji funkcji Greena
są więc spełnione w zmiennej $x$.
Ponieważ $Q_R(x,y)=Q_R(y,x)$, mamy
$H_R(x,y)=H_R(y,x)$, a radialność daje
$\Phi_n(x-y)=\Phi_n(y-x)$. Stąd
$G_R(x,y)=G_R(y,x)$, także wtedy, gdy jeden z punktów
jest zerem. Symetria pozwala stosować tę samą konstrukcję
w zmiennej $y$ przy ustalonym $x$; harmoniczna poprawka
$H_R(x,\cdot)$ jest gładka na otoczeniu domkniętej kuli.
W szczególności spełnione są założenia dotyczące
funkcji Greena we wzorze reprezentacyjnym
\eqref{eq:pde3-green-representation}.
\end{proof}

\begin{center}
\begin{tikzpicture}[scale=.87,>=Stealth]
 \fill[blue!5] (0,0) circle (2);
 \draw[thick] (0,0) circle (2);
 \draw[densely dashed,gray!65] (0,0)--(4.25,0);
 \fill (0,0) circle (1.6pt);
 \node[below left] at (0,0) {$0$};
 \fill[blue!75!black] (1.1,0) circle (2.5pt);
 \node[above] at (1.1,0) {$y$};
 \fill[red!75!black] (3.63,0) circle (2.5pt);
 \node[above] at (3.63,0) {$y^*$};
 \draw[<->,thin] (0,-.3)--(2,-.3);
 \node[below] at (1,-.3) {$R$};
 \node[align=left,anchor=west] at (4.2,1.0) {$|y|\,|y^*|=R^2$\\punkt odwrócony leży poza kulą};
\end{tikzpicture}
\par\smallskip\small
Inwersja w sferze $\partial B_R$: punkt $y$ wewnątrz
odpowiada punktowi $y^*$ na zewnątrz. Pozwala to
zbudować harmoniczną poprawkę rozwiązania fundamentalnego.
\end{center}

\begin{twierdzenie}[Jądro Poissona dla kuli]
\label{thm:pde3-poisson}
Niech $n\geq2$, $R>0$ i $g\in C(\partial B_R)$.
Dla $x\in B_R$ oraz $\zeta\in\partial B_R$ zdefiniujmy
jądro $P_R$ i funkcję $u$ na $B_R$:
\begin{equation}\label{eq:pde3-poisson}
 P_R(x,\zeta)=\frac{R^2-|x|^2}{\omega_n R|x-\zeta|^n},
 \qquad
 u(x)=\int_{\partial B_R}P_R(x,\zeta)g(\zeta)
 \,\mathrm d\sigma(\zeta).
\end{equation}
Wówczas $u\in C^\infty(B_R)\cap C(\overline{B_R})$,
$\Delta u=0$ w $B_R$ i $u|_{\partial B_R}=g$.
Jest to jedyne rozwiązanie z klasy
$C^2(B_R)\cap C(\overline{B_R})$ tego problemu Dirichleta.
Ponadto $P_R(x,\zeta)>0$ oraz
$\int_{\partial B_R}P_R(x,\zeta)\,\mathrm d\sigma(\zeta)=1$.
\end{twierdzenie}
\begin{proof}
Najpierw wykażemy związek z normalną pochodną funkcji
Greena. Ustalmy biegun $y\in B_R$ oraz punkt
$\zeta\in\partial B_R$. Jednostkowa normalna zewnętrzna
do kuli wynosi $\nu_\zeta=\zeta/R$.
Korzystamy z jednolitej postaci poprawki harmonicznej
z dowodu stwierdzenia~\ref{prop:pde3-ball-green}:
\[
 G_R(x,y)=\Phi_n(x-y)-h_n(\sqrt{Q_R(x,y)}).
\]
W otoczeniu $\zeta$ obie funkcje po prawej stronie
są gładkie w zmiennej $x$, ponieważ $\zeta\ne y$
i $Q_R(\zeta,y)=|\zeta-y|^2>0$.
Wzór $h_n'(r)=-1/(\omega_n r^{n-1})$ jest prawdziwy
zarówno dla $n=2$, jak i $n\geq3$.
Ponadto
$\nabla_x Q_R(x,y)=-2y+2|y|^2x/R^2$.
Z reguły łańcuchowej otrzymujemy
\[
 \nabla_x H_R(x,y)=
 -\frac{|y|^2x/R^2-y}{\omega_n Q_R(x,y)^{n/2}},
 \qquad H_R(x,y)=h_n(\sqrt{Q_R(x,y)}).
\]
Wobec $\nabla\Phi_n(z)=-z/(\omega_n|z|^n)$
oraz $Q_R(\zeta,y)=|\zeta-y|^2$ mamy
\[
 \begin{aligned}
 -\nabla_xG_R(\zeta,y)\cdot\nu_\zeta
 &=\frac{(\zeta-y)\cdot\zeta
       -( |y|^2\zeta/R^2-y)\cdot\zeta}
       {\omega_n R|\zeta-y|^n}\\
 &=\frac{R^2-|y|^2}{\omega_n R|\zeta-y|^n}.
 \end{aligned}
\]
Rachunek obejmuje również $n=2$ oraz $y=0$,
bez przechodzenia do granicy w formule inwersyjnej.
Dla bieguna $y=0$ poprawka $H_R(x,0)$ jest stała,
a pochodna normalna samego rozwiązania fundamentalnego
daje $1/(\omega_n R^{n-1})$ po zmianie znaku.
Z symetrii funkcji Greena otrzymujemy
\[
 P_R(x,\zeta)=-\nabla_yG_R(x,\zeta)\cdot\nu_\zeta,
 \qquad x\in B_R,\quad \zeta\in\partial B_R,
\]
gdzie pochodna dotyczy drugiej zmiennej i jest obliczona
przez gładkie przedłużenie w otoczeniu brzegu.

Wykażmy teraz harmoniczność jądra względem $x$.
Ustalmy niepusty zwarty zbiór $K\subset B_R$ i
$\zeta\in\partial B_R$. Istnieje $\rho<R$ taka, że
$K\subset B_\rho$. Wybierzmy $\eta>0$ tak małe, aby
$\eta<R-\rho$ i $\rho(R+\eta)<R^2$.
Dla $x\in B_\rho$ i $y\in B_\eta(\zeta)$ mamy
$|x-y|\geq R-\eta-\rho>0$ i $|x||y|<R^2$.
Jednolity wzór na $G_R$ jest więc gładki na iloczynie
tych otwartych zbiorów. Dla $y\in B_\eta(\zeta)\cap B_R$
mamy $\Delta_xG_R(x,y)=0$ w $B_\rho$.
Różniczkując tę tożsamość względem $y$ w ustalonym
kierunku $\nu_\zeta$ i przechodząc do $y=\zeta$
z wnętrza kuli, dostajemy
\[
 \Delta_xP_R(x,\zeta)=
 -\nu_\zeta\cdot\nabla_y\Delta_xG_R(x,\zeta)=0.
\]
Komutacja pochodnych i przejście graniczne są
uzasadnione gładkością na wskazanym iloczynie.
Nie różniczkujemy tu ograniczenia $G_R(x,\cdot)$
do samej sfery, lecz funkcję określoną także
w kierunku normalnym.

Dla każdego wielowskaźnika $\alpha\in\mathbb N_0^n$
pochodna $D_x^\alpha P_R$ jest ciągła i ograniczona
na $\overline{B_\rho}\times\partial B_R$, ponieważ
$|x-\zeta|\geq R-\rho>0$. Funkcja $g$ jest ograniczona,
a sfera ma skończoną miarę. Zatem różniczkowanie
całki definiującej $u$ względem $x$ dowolną liczbę razy
jest dopuszczalne na $B_\rho$ na mocy twierdzenia
o zbieżności zdominowanej, stosowanego również
do ilorazów różnicowych. W szczególności
$\Delta u(x)=\int_{\partial B_R}\Delta_x P_R(x,\zeta)
g(\zeta)\,\mathrm d\sigma(\zeta)=0$.
Dowolność $\rho<R$ daje $u\in C^\infty(B_R)$
i harmoniczność w całej kuli.

Uzasadnijmy normalizację jądra bez powoływania się na
istnienie szukanego rozwiązania Dirichleta. Stosujemy
wzór reprezentacyjny~\eqref{eq:pde3-green-representation}
do funkcji stałej $1$, dla której $-\Delta1=0$.
Poprawka harmoniczna w funkcji Greena dla kuli
jest gładka na otoczeniu domkniętej kuli przy ustalonym
biegunie $x$, więc założenia wzoru są spełnione.
Wynika stąd $1=\int_{\partial B_R}P_R(x,\zeta)\,
\mathrm d\sigma(\zeta)$ dla każdego $x\in B_R$.
Dodatniość wynika z $R^2-|x|^2>0$, dodatniości
$\omega_n R$ oraz $|x-\zeta|>0$.
W szczególności całka Poissona jest średnią ważoną
danych brzegowych i $|u(x)|\leq\|g\|_\infty$.

Ustalmy $\zeta_0\in\partial B_R$ i $\eps>0$.
Z jednostajnej ciągłości $g$ wybierzmy $\delta>0$
tak, że $|g(\zeta)-g(\zeta_0)|<\eps$, gdy
$|\zeta-\zeta_0|<\delta$. Posługując się dodatniością
i normalizacją $P_R$ dostajemy
\[
 |u(x)-g(\zeta_0)|\leq\eps+
 2\|g\|_\infty
 \int_{\substack{\zeta\in\partial B_R\\
 |\zeta-\zeta_0|\geq\delta}}P_R(x,\zeta)\,
 \mathrm d\sigma(\zeta).
\]
Jeżeli $|x-\zeta_0|<\delta/2$, to na zbiorze
całkowania $|x-\zeta|\geq|\zeta-\zeta_0|
-|x-\zeta_0|\geq\delta/2$.
Zatem
\[
 \int_{|\zeta-\zeta_0|\geq\delta}P_R(x,\zeta)\,
 \mathrm d\sigma(\zeta)
 \leq \frac{(R^2-|x|^2)\omega_n R^{n-1}}
 {\omega_n R(\delta/2)^n}\longrightarrow0
 \quad(x\to\zeta_0,\ x\in B_R),
\]
gdzie całkujemy po odpowiedniej części $\partial B_R$.
Stąd $\limsup_{x\to\zeta_0,\ x\in B_R}
|u(x)-g(\zeta_0)|\leq\eps$. Dowolność $\eps>0$
daje żądaną granicę. Definiując $u=g$ na brzegu,
otrzymujemy funkcję ciągłą na $\overline{B_R}$,
gdyż $g$ jest ciągła na sferze, a granice od wnętrza
zgadzają się z jej wartościami.
Jeżeli $v$ jest innym rozwiązaniem w podanej klasie,
to $u-v$ jest harmoniczna w kuli i znika na brzegu.
Słaba zasada maksimum (twierdzenie~\ref{thm:laplace-max})
zastosowana do $u-v$ oraz $v-u$ daje $u=v$.
\end{proof}

\begin{przyklad}[Równanie Dirichleta na dysku]
Niech $n=2$, $R=1$, a $g(\cos\theta,\sin\theta)=\cos(3\theta)$.
Połóżmy $u(x_1,x_2)=x_1^3-3x_1x_2^2$.
Jest to wielomian, więc jest gładki także w środku dysku.
Mamy $\partial_{x_1}^2u=6x_1$ oraz
$\partial_{x_2}^2u=-6x_1$, skąd $\Delta u=0$.
Tożsamość $\cos(3\theta)=\cos^3\theta
-3\cos\theta\sin^2\theta$ daje
$u(r\cos\theta,r\sin\theta)=r^3\cos(3\theta)$.
Dla $r=1$ otrzymujemy dane $g$.
Z jednoznaczności w twierdzeniu~\ref{thm:pde3-poisson}
funkcja ta jest dokładnie całką Poissona z tych danych.
\end{przyklad}

\begin{stwierdzenie}[Nierówność Harnacka w kuli]
\label{prop:pde3-harnack}
Niech $n\geq2$, $R>0$, $u\in C^2(B_R)\cap C(\overline{B_R})$
będzie harmoniczna i nieujemna. Dla każdego $0\leq r<R$
i $|x|\leq r$ zachodzi
\begin{equation}\label{eq:pde3-harnack}
 \frac{R^{n-2}(R-r)}{(R+r)^{n-1}}u(0)
 \leq u(x)\leq
 \frac{R^{n-2}(R+r)}{(R-r)^{n-1}}u(0).
\end{equation}
\end{stwierdzenie}
\begin{proof}
Ciągłość $u$ na domkniętej kuli i nieujemność we wnętrzu
dają $u|_{\partial B_R}\geq0$. Twierdzenie
\ref{thm:pde3-poisson} i jednoznaczność rozwiązania
zapewniają reprezentację $u$ jako całki Poissona
z własnych wartości brzegowych. Możemy więc porównać
jądro w punkcie $x$ i w punkcie $0$.
Ponieważ $P_R(0,\zeta)=1/(\omega_n R^{n-1})$,
\[
 \frac{P_R(x,\zeta)}{P_R(0,\zeta)}
 =\frac{R^{n-2}(R^2-|x|^2)}{|x-\zeta|^n}.
\]
Używając $R-|x|\leq|x-\zeta|\leq R+|x|$,
otrzymujemy
\[
 \frac{R^{n-2}(R-|x|)}{(R+|x|)^{n-1}}
 \leq\frac{P_R(x,\zeta)}{P_R(0,\zeta)}
 \leq\frac{R^{n-2}(R+|x|)}{(R-|x|)^{n-1}}.
\]
Dla $0\leq s<R$ funkcja
$(R-s)/(R+s)^{n-1}$ jest malejąca, ponieważ jej
dodatni licznik maleje, a dodatni mianownik rośnie.
Analogicznie $(R+s)/(R-s)^{n-1}$ jest rosnąca.
Podstawiając $s=|x|\leq r$, otrzymujemy oszacowania
niezależne od $\zeta$. Mnożymy nierówności przez
nieujemne $P_R(0,\zeta)u(\zeta)$ i całkujemy po
$\partial B_R$. Całka z $P_R(0,\zeta)u(\zeta)$
jest równa $u(0)$, a z $P_R(x,\zeta)u(\zeta)$
jest równa $u(x)$. Otrzymujemy
\eqref{eq:pde3-harnack}, również wtedy, gdy $u(0)=0$,
bez dzielenia przez wartość funkcji.
\end{proof}

\begin{uwaga}[Nierówność Harnacka dla kul wewnętrznych]
Niech $\Omega\subset\R^n$ będzie obszarem, a
$u\in C^2(\Omega)$ nieujemną funkcją harmoniczną.
Nierówność Harnacka obowiązuje po translacji
na każdej kuli o domknięciu zawartym w $\Omega$.
W szczególności dla $\overline{B_{2a}(z)}\subset\Omega$
i $w\in\overline{B_a(z)}$ wstawienie $R=2a$, $r=a$
do \eqref{eq:pde3-harnack} daje
\[
 c_n u(z)\leq u(w)\leq C_n u(z),\qquad
 c_n=\frac{2^{n-2}}{3^{n-1}},\quad C_n=3\cdot2^{n-2}.
\]
Wykażmy, że dla każdego niepustego zwartego
$K\subset\Omega$ istnieje stała $C_K\geq1$,
niezależna od $u$, taka, że
$\sup_K u\leq C_K\inf_K u$.

Ustalmy $z_0\in\Omega$. Dla każdego $z\in K$
istnieje droga ciągła $\gamma\colon[0,1]\to\Omega$
od $z_0$ do $z$, ponieważ obszar w $\R^n$ jest
spójny drogowo. Jej obraz jest zwarty; można więc
wybrać $a>0$ tak, aby
$\overline{B_{2a}(\gamma(t))}\subset\Omega$ dla
wszystkich $t\in[0,1]$.
Z jednostajnej ciągłości drogi istnieje podział
$0=t_0<\cdots<t_m=1$, dla którego
$|\gamma(t_{j+1})-\gamma(t_j)|<a$.
Wzdłuż tych punktów kolejne zastosowania oszacowania
dają $c_n^m u(z_0)\leq u(z)\leq C_n^m u(z_0)$.
Jedno dalsze zastosowanie daje dla $w\in B_a(z)$
nierówności
$c_n^{m+1}u(z_0)\leq u(w)\leq C_n^{m+1}u(z_0)$.

Kule $B_a(z)$ uzyskane w ten sposób pokrywają $K$;
ze zwartości wybieramy skończone podpokrycie.
Niech $N$ będzie największą liczbą $m+1$
wśród wybranych łańcuchów. Ponieważ
$0<c_n\leq1\leq C_n$, dla wszystkich $w\in K$
mamy $c_n^N u(z_0)\leq u(w)\leq C_n^N u(z_0)$.
Ostatecznie $\sup_Ku\leq (C_n/c_n)^N\inf_Ku$.
Konstrukcja stałej zależy wyłącznie od geometrii
wybranych kul i dróg, a nie od funkcji $u$.
Nie zakładamy ciągłego przedłużenia $u$ na
\emph{brzeg całego obszaru}.
\end{uwaga}

\section{Problem Poissona i bariery porównawcze}
\label{sec:pde3-poisson-equation}

\begin{twierdzenie}[Problem Poissona na kuli]
\label{thm:pde3-ball-poisson}
Niech $n\geq2$, $R>0$, $f\in C_c^\infty(B_R)$,
a $g\in C(\partial B_R)$. Istnieje dokładnie jedno
$u\in C^2(B_R)\cap C(\overline{B_R})$ takie, że
$-\Delta u=f$ w $B_R$ i $u=g$ na $\partial B_R$.
Ponadto $u\in C^\infty(B_R)$ i dla każdego $x\in B_R$
ma postać
\begin{equation}\label{eq:pde3-poisson-solution}
 u(x)=\int_{B_R}G_R(x,y)f(y)\,\mathrm d\lambda_n(y)
 +\int_{\partial B_R}P_R(x,\zeta)g(\zeta)
 \,\mathrm d\sigma(\zeta).
\end{equation}
\end{twierdzenie}
\begin{proof}
Niech $\widetilde f\colon\R^n\to\R$ oznacza przedłużenie
$f$ przez zero. Ponieważ $\supp f$ jest zwartym
podzbiorem $B_R$, funkcja $f$ znika w pewnym otoczeniu
$\partial B_R$ po wewnętrznej stronie sfery.
Stąd $\widetilde f\in C_c^\infty(\R^n)$.
Połóżmy $w=\Phi_n*\widetilde f$. Z twierdzenia
\ref{thm:pde3-newton} wynika, że $w\in C^\infty(\R^n)$
oraz $-\Delta w=\widetilde f$ na $\R^n$.

Dane $g-w|_{\partial B_R}$ są ciągłe na sferze.
Twierdzenie~\ref{thm:pde3-poisson} dostarcza funkcji
$h\in C^\infty(B_R)\cap C(\overline{B_R})$ harmonicznej
w kuli i równej $g-w$ na jej brzegu. Funkcja $u=w+h$
spełnia równanie $-\Delta u=f$ i warunek brzegowy $u=g$.
Jeśli $u_1,u_2\in C^2(B_R)\cap C(\overline{B_R})$
są dwoma rozwiązaniami, ich różnica jest harmoniczna
i znika na brzegu. Słaba zasada maksimum zastosowana
do tej różnicy i jej przeciwności daje $u_1=u_2$.

Wykażmy wzór~\eqref{eq:pde3-poisson-solution}.
Korzystamy z poprawki harmonicznej
$H_R(x,y)=h_n(\sqrt{Q_R(x,y)})$ skonstruowanej
w dowodzie stwierdzenia~\ref{prop:pde3-ball-green},
gdzie $Q_R(x,y)=R^2-2x\cdot y+|x|^2|y|^2/R^2$.
Dla $x\ne y$ mamy $G_R(x,y)=\Phi_n(x-y)-H_R(x,y)$.
Wybierzmy $0<r<R$ tak, aby $\supp\widetilde f\subset B_r$
(gdy $f=0$, można wybrać dowolne takie $r$), oraz
$\eta>0$ spełniające $(R+\eta)r<R^2$. Dla
$|x|\leq R+\eta$ i $|y|\leq r$ zachodzi
\[
 Q_R(x,y)\geq
 \left(R-\frac{(R+\eta)r}{R}\right)^2>0.
\]
Poprawka $H_R$ jest więc gładka na otoczeniu zwartego
zbioru $\overline{B_{R+\eta}}\times\overline{B_r}$.
Dla każdego wielowskaźnika $\alpha\in\mathbb N_0^n$
istnieje skończona stała
\[
 M_\alpha=\sup_{\substack{|x|\leq R+\eta\\ |y|\leq r}}
 |D_x^\alpha H_R(x,y)|.
\]
Definiujemy dla $x\in B_{R+\eta}$
\[
 h_f(x)=\int_{B_r}H_R(x,y)f(y)\,\mathrm d\lambda_n(y).
\]
Funkcja $M_\alpha|f(y)|$ jest całkowalna w $B_r$
i dominuje moduł odpowiedniej pochodnej funkcji
podcałkowej. Twierdzenie o zbieżności zdominowanej,
stosowane do ilorazów różnicowych i do samych pochodnych,
pozwala różniczkować pod całką dowolną liczbę razy.
W konsekwencji $h_f\in C^\infty(B_{R+\eta})$ oraz
\[
 D^\alpha h_f(x)=\int_{B_r}D_x^\alpha H_R(x,y)
 f(y)\,\mathrm d\lambda_n(y).
\]
Dla każdego $y\in B_r$ funkcja $H_R(\cdot,y)$
jest harmoniczna w $B_R$, również dla $y=0$.
Stąd $\Delta h_f=0$ w $B_R$.

Dla $x\in B_R$ obie funkcje podcałkowe
$\Phi_n(x-y)f(y)$ i $H_R(x,y)f(y)$ są całkowalne:
pierwsza dzięki lokalnej całkowalności rozwiązania
fundamentalnego i ograniczonemu nośnikowi danych,
druga dzięki powyższemu oszacowaniu z $\alpha=0$.
Punkt $y=x$ ma miarę Lebesgue'a równą zeru, więc
dowolna wartość funkcji Greena w tym punkcie nie
zmienia całki. Zatem funkcja
\[
 v(x)=w(x)-h_f(x)=
 \int_{B_R}G_R(x,y)f(y)\,\mathrm d\lambda_n(y)
\]
jest gładka w kuli i spełnia $-\Delta v=f$.
Ponadto $w$ i $h_f$ są ciągłe na domkniętej kuli.
Jeśli $x\in\partial B_R$, to $|x-y|\geq R-r>0$
dla $y\in\supp f$, a $H_R(x,y)=\Phi_n(x-y)$.
Wynika stąd bezpośrednio, że $w(x)=h_f(x)$,
czyli $v(x)=0$ na brzegu.

Niech $p_g$ będzie całką Poissona z danych $g$.
Funkcja $v+p_g$ jest gładka w kuli, ciągła na jej
domknięciu, spełnia $-\Delta(v+p_g)=f$ i jest równa
$g$ na brzegu. Z wykazanej już jednoznaczności
$u=v+p_g$, co daje~\eqref{eq:pde3-poisson-solution}.
\end{proof}

\begin{stwierdzenie}[Oszacowanie przez barierę kwadratową]
\label{prop:pde3-barrier}
Niech $n\geq2$, $R>0$ i niech $\Omega\subset B_R(0)$
będzie obszarem. Niech $f\in C(\Omega)$ będzie rzeczywista
i ograniczona, a $u\in C^2(\Omega)\cap C(\overline\Omega)$
będzie rzeczywistym rozwiązaniem równania $-\Delta u=f$
w $\Omega$. Wówczas
\begin{equation}\label{eq:pde3-barrier}
 \sup_{x\in\Omega}|u(x)|\leq
 \sup_{x\in\partial\Omega}|u(x)|+
 \frac{R^2}{2n}\sup_{x\in\Omega}|f(x)|.
\end{equation}
\end{stwierdzenie}
\begin{proof}
Oznaczmy $M=\sup_\Omega|f|$ oraz $A=\sup_{\partial\Omega}|u|$.
Liczba $M$ jest skończona z założenia, a $A$ jest
skończona dzięki ciągłości $u$ na zwartym zbiorze
$\overline\Omega$. Połóżmy $b(x)=(R^2-|x|^2)/(2n)$
dla $x\in\R^n$. Ponieważ
$\overline\Omega\subset\overline{B_R(0)}$,
mamy $0\leq b\leq R^2/(2n)$ na $\overline\Omega$.
Dla każdego $j=1,\ldots,n$ zachodzi
$\partial_{x_j}^2b=-1/n$, zatem $-\Delta b=1$.

Funkcja $w_+=u-Mb$ należy do
$C^2(\Omega)\cap C(\overline\Omega)$ i spełnia
$\Delta w_+=-f+M\geq0$. Na brzegu mamy
$w_+\leq u\leq A$, gdyż $M\geq0$ i $b\geq0$.
Ze słabej zasady maksimum, twierdzenie~\ref{thm:laplace-max},
otrzymujemy $w_+\leq A$ w $\Omega$.
Analogicznie $w_-=-u-Mb$ spełnia
$\Delta w_-=f+M\geq0$ oraz $w_-\leq A$ na brzegu.
Stąd $w_-\leq A$ w $\Omega$.
Obie nierówności dają punktowe oszacowanie
\[
 |u(x)|\leq A+Mb(x)
 =A+\frac{R^2-|x|^2}{2n}M
 \qquad(x\in\Omega).
\]
Biorąc supremum po $x\in\Omega$ i używając
$b(x)\leq R^2/(2n)$, dostajemy
\eqref{eq:pde3-barrier}. Nie wymagamy przy tym
gładkości brzegu obszaru.
\end{proof}

\begin{przyklad}[Wymuszenie stałe]
Niech $n\geq2$, $R>0$, $f\equiv1$ w $B_R(0)$,
a dane Dirichleta będą równe $0$.
Wtedy $u(x)=(R^2-|x|^2)/(2n)$ jest wielomianem
spełniającym $-\Delta u=1$ w kuli i $u=0$ na sferze.
Z zasady maksimum jest to jedyne rozwiązanie
w klasie $C^2(B_R)\cap C(\overline{B_R})$.
Mamy $\sup_{B_R}|u|=u(0)=R^2/(2n)$,
$\sup_{\partial B_R}|u|=0$ i $\sup_{B_R}|f|=1$.
Zatem w nierówności~\eqref{eq:pde3-barrier} zachodzi
równość; współczynnika $R^2/(2n)$ przy normie wymuszenia
nie można zmniejszyć w całej rozważanej klasie
obszarów i danych.
Zauważmy, że $f\equiv1$ w kuli nie ma w niej zwartego
nośnika, więc istnienia rozwiązania w tym przykładzie
nie wywodzimy z twierdzenia~\ref{thm:pde3-ball-poisson},
lecz sprawdzamy je bezpośrednio.
\end{przyklad}

\section{Harmoniczność, analityczność i funkcje holomorficzne}
\label{sec:pde3-analytic}

\begin{twierdzenie}[Analityczność funkcji harmonicznych]
\label{thm:pde3-analytic}
Niech $n\geq2$, $\Omega\subset\R^n$ będzie otwarta,
a $u\in C^2(\Omega)$ harmoniczna. Wtedy $u$ jest
rzeczywiście analityczna w $\Omega$, czyli w otoczeniu
każdego punktu jest sumą zbieżnego szeregu potęgowego.
W szczególności $u\in C^\infty(\Omega)$.
\end{twierdzenie}
\begin{proof}
Ustalmy $x_0\in\Omega$ i wybierzmy $R>0$ tak, aby
$\overline{B_R(x_0)}\subset\Omega$.
Z twierdzenia~\ref{thm:pde3-poisson}, zastosowanego
po translacji do kuli o środku $x_0$, i z jednoznaczności
rozwiązania zagadnienia Dirichleta mamy
\[
 u(x)=\int_{\partial B_R(x_0)}
 \frac{R^2-|x-x_0|^2}{\omega_n R|x-\zeta|^n}
 u(\zeta)\,\mathrm d\sigma(\zeta)
 \qquad(x\in B_R(x_0)).
\]
Oznaczmy $h=x-x_0$ i $\theta=(\zeta-x_0)/R\in S^{n-1}$.
Po zmianie zmiennych na sferze otrzymujemy
\[
 u(x_0+h)=\frac{1-|h|^2/R^2}{\omega_n}
 \int_{S^{n-1}}(1+q_\theta(h))^{-n/2}
 u(x_0+R\theta)\,\mathrm d\sigma(\theta),
\]
gdzie $q_\theta(h)=-2\theta\cdot h/R+|h|^2/R^2$.
Dla ustalonego $\theta$ i $|q_\theta(h)|<1$ stosujemy
rozwinięcie dwumianowe
\[
 (1+q_\theta(h))^{-n/2}
 =\sum_{k=0}^\infty a_k q_\theta(h)^k,\qquad
 a_0=1,\quad
 a_k=\frac{(-1)^k}{k!}\prod_{\ell=0}^{k-1}
 \left(\frac n2+\ell\right)\quad(k\geq1).
\]

Aby uzasadnić, że otrzymujemy rzeczywiście szereg
potęgowy w $n$ zmiennych, a nie tylko jednostajnie
zbieżny szereg wielomianów, oszacujemy również
sumę modułów ich współczynników z wagami.
Dla wielomianu $p(h)=\sum_\alpha p_\alpha h^\alpha$
i $\rho>0$ oznaczmy
\[
 L_\rho(p)=\sum_\alpha |p_\alpha|\rho^{|\alpha|},
 \qquad h^\alpha=\prod_{j=1}^n h_j^{\alpha_j},
 \quad |\alpha|=\sum_{j=1}^n\alpha_j,
\]
gdzie $\alpha\in\mathbb N_0^n$.
Z reguły mnożenia wielomianów i nierówności trójkąta
wynika $L_\rho(p_1p_2)\leq L_\rho(p_1)L_\rho(p_2)$.
Ponadto $|p(h)|\leq L_\rho(p)$, gdy $|h_j|\leq\rho$
dla każdego $j$. Wybierzmy $\rho=R/(4n)$.
Ponieważ $|\theta_j|\leq1$, mamy niezależnie od $\theta$
\[
 L_\rho(q_\theta)
 \leq\frac{2n\rho}{R}+\frac{n\rho^2}{R^2}
 =\frac12+\frac1{16n}\leq\frac9{16}=:c<1.
\]
Domknięty sześcian $|h_j|\leq\rho$ zawiera się
w $B_R(0)$, gdyż $|h|\leq\sqrt n\,\rho<R$.
Na tym sześcianie rozwinięcie dwumianowe jest poprawne.
Co więcej,
\[
 \sum_{k=0}^\infty |a_k|L_\rho(q_\theta^k)
 \leq\sum_{k=0}^\infty |a_k|c^k<\infty,
\]
ponieważ iloraz kolejnych wyrazów szeregu po prawej
dąży do $c<1$.

Niech $b_{k,\alpha}(\theta)$ oznacza współczynnik
przy $h^\alpha$ w wielomianie $q_\theta^k$.
Są to ciągłe funkcje $\theta$, a dla ustalonego
$\alpha$ jedynie skończenie wiele $k$ może dać
niezerowy współczynnik, ponieważ każdy jednomian
w $q_\theta^k$ ma stopień co najmniej $k$.
Połóżmy
\[
 d_\alpha=\frac1{\omega_n}
 \sum_{k=0}^{|\alpha|}a_k
 \int_{S^{n-1}} b_{k,\alpha}(\theta)u(x_0+R\theta)
 \,\mathrm d\sigma(\theta).
\]
Dla $M=\sup_{\partial B_R(x_0)}|u|<\infty$
nierówność trójkąta i twierdzenie Tonellego dla
nieujemnych składników dają
\[
 \sum_\alpha |d_\alpha|\rho^{|\alpha|}
 \leq M\sum_{k=0}^\infty |a_k|
 \sup_{\theta\in S^{n-1}}L_\rho(q_\theta^k)
 \leq M\sum_{k=0}^\infty |a_k|c^k<\infty.
\]
To samo oszacowanie uzasadnia całkowanie wyraz po
wyrazie i zmianę kolejności sumowania rozwinięcia
dwumianowego oraz jednomianów. Otrzymujemy
\[
 u(x_0+h)=\left(1-\frac{|h|^2}{R^2}\right)
 \sum_{\alpha\in\mathbb N_0^n}d_\alpha h^\alpha
 \qquad(|h_j|\leq\rho,\ j=1,\ldots,n).
\]
Mnożenie przez skończony wielomian zachowuje
bezwzględną zbieżność z wagami $\rho^{|\alpha|}$
na mocy tej samej nierówności dla współczynników.
W szczególności otrzymujemy szereg potęgowy
przedstawiający $u$ w otwartym otoczeniu $x_0$,
zbieżny bezwzględnie i jednostajnie na wskazanym
domkniętym sześcianie. Na każdym mniejszym sześcianie
można różniczkować ten szereg wyraz po wyrazie
dowolną liczbę razy: czynniki wielomianowe w
$|\alpha|$ pochodzące od różniczkowania są
kontrolowane przez geometryczny czynnik
$(\rho'/\rho)^{|\alpha|}$ dla $0<\rho'<\rho$.
Otrzymane współczynniki są zatem współczynnikami
Taylora funkcji $u$ w $x_0$.
Dowolność $x_0$ daje analityczność i gładkość $u$.
\end{proof}

\begin{stwierdzenie}[Harmoniczna sprzężona w wymiarze dwa]
\label{prop:pde3-conjugate}
Niech $\Omega\subset\R^2$ będzie obszarem jednospójnym,
a $u\in C^2(\Omega)$ rzeczywistą funkcją harmoniczną.
Istnieje rzeczywista funkcja harmoniczna $v\in C^2(\Omega)$,
jednoznaczna z dokładnością do stałej addytywnej, taka,
że $F(x+iy)=u(x,y)+iv(x,y)$ jest holomorficzna w $\Omega$
traktowanej jako podzbiór $\C$.
\end{stwierdzenie}
\begin{proof}
Rozważmy formę różniczkową
$\alpha=-u_y\,\mathrm dx+u_x\,\mathrm dy$ klasy $C^1$.
Korzystając z reguły różniczkowania zewnętrznego
i antysymetrii iloczynu zewnętrznego, obliczamy
\[
 \mathrm d\alpha
 =-(u_{yx}\,\mathrm dx+u_{yy}\,\mathrm dy)
       \wedge\mathrm dx
 +(u_{xx}\,\mathrm dx+u_{xy}\,\mathrm dy)
       \wedge\mathrm dy
 =\Delta u\,\mathrm dx\wedge\mathrm dy=0.
\]
Stosujemy twierdzenie o dokładności zamkniętych form
różniczkowych stopnia $1$ klasy $C^1$ na jednospójnym
obszarze: całka takiej formy wzdłuż drogi kawałkami
klasy $C^1$ zależy jedynie od jej końców.
Ustalmy $(x_*,y_*)\in\Omega$ i $c\in\R$.
Definiujemy $v(x,y)=c+\int_\gamma\alpha$, gdzie
$\gamma$ jest dowolną drogą kawałkami klasy $C^1$
w $\Omega$ od $(x_*,y_*)$ do $(x,y)$.
Jednospójność i zamkniętość formy zapewniają
niezależność definicji od wyboru drogi.

W otoczeniu dowolnego punktu $(x,y)\in\Omega$
możemy do wybranej drogi dołączyć krótki odcinek
równoległy do jednej z osi, zawarty w $\Omega$.
Z addytywności całki krzywoliniowej otrzymujemy
\[
 \begin{aligned}
 v(x+h,y)-v(x,y)&=-\int_0^h u_y(x+t,y)\,
                       \mathrm d\lambda_1(t),\\
 v(x,y+h)-v(x,y)&=\int_0^h u_x(x,y+t)\,
                       \mathrm d\lambda_1(t).
 \end{aligned}
\]
Dla $h<0$ całki rozumiemy z odwróconą orientacją.
Ciągłość pochodnych $u$ daje po podzieleniu przez $h$
i przejściu do granicy
$v_x=-u_y$ i $v_y=u_x$. Obie pochodne są klasy $C^1$,
zatem $v\in C^2(\Omega)$ i $\mathrm dv=\alpha$.
Równość pochodnych mieszanych $u$ daje
$\Delta v=-u_{yx}+u_{xy}=0$.

Równości $u_x=v_y$ i $u_y=-v_x$ są równaniami
Cauchy'ego--Riemanna. Przy regularności $C^1$
funkcji $F$ zapewniają jej holomorficzność.
Jeżeli $v_1,v_2$ są dwiema funkcjami sprzężonymi
do $u$, to z równań Cauchy'ego--Riemanna wynika
$\nabla(v_1-v_2)=0$. Różnica jest lokalnie stała,
a ponieważ $\Omega$ jest spójna, jest stała na
całym obszarze. Odwrotnie, dodanie dowolnej stałej
rzeczywistej do $v$ zachowuje wszystkie warunki.
\end{proof}

\begin{kontrprzyklad}[Przeszkoda topologiczna dla sprzężenia]
Na pierścieniu $\Omega=\{z\in\C:1<|z|<2\}$
połóżmy $u(x,y)=\frac12\log(x^2+y^2)$.
Oznaczając $s=x^2+y^2$, obliczamy
\[
 u_x=\frac{x}{s},\quad u_y=\frac{y}{s},\qquad
 u_{xx}=\frac{y^2-x^2}{s^2},\quad
 u_{yy}=\frac{x^2-y^2}{s^2}.
\]
Stąd $\Delta u=0$ na $\Omega$. Gdyby istniała
globalna funkcja sprzężona $v$ do $u$, równania
Cauchy'ego--Riemanna wymuszałyby
\[
 \mathrm dv=\alpha
 =-\frac{y}{x^2+y^2}\,\mathrm dx
 +\frac{x}{x^2+y^2}\,\mathrm dy.
\]
Niech $a=3/2$ i $\gamma(t)=(a\cos t,a\sin t)$
dla $0\leq t\leq2\pi$. Jest to zamknięta krzywa
w $\Omega$, zorientowana przeciwnie do ruchu
wskazówek zegara. Mamy
\[
 \int_\gamma\alpha
 =\int_0^{2\pi}
 \frac{a^2\sin^2t+a^2\cos^2t}{a^2}
 \,\mathrm d\lambda_1(t)=2\pi.
\]
Tymczasem dla formy dokładnej $\mathrm dv$
całka wzdłuż $\gamma$ byłaby równa
$v(\gamma(2\pi))-v(\gamma(0))=0$.
Sprzeczność dowodzi, że nie istnieje globalna
jednowartościowa funkcja sprzężona do $u$.
Każdy punkt pierścienia ma jednak otoczenie
będące kulą o domknięciu zawartym w $\Omega$;
na takiej kuli istnieje funkcja sprzężona na mocy
stwierdzenia~\ref{prop:pde3-conjugate}.
\end{kontrprzyklad}

\section{Wygładzanie harmoniczne na dysku: rozwiązanie przez szeregi Fouriera}
\label{sec:pde3-fourier-disk}

Jawna postać jądra Poissona ma szczególnie przejrzystą
interpretację na jednostkowym dysku. Niech
$g\in C^2(S^1)$ będzie funkcją rzeczywistą; oznacza to,
że funkcja $G(\theta)=g(e^{i\theta})$ jest klasy $C^2$
i $2\pi$-okresowa na $\R$. Ustalamy normalizację
współczynników Fouriera:
\[
 \begin{aligned}
 a_0&=\frac1{2\pi}\int_0^{2\pi}G(\theta)\,
             \mathrm d\lambda_1(\theta),\\
 a_k&=\frac1\pi\int_0^{2\pi}G(\theta)\cos(k\theta)\,
             \mathrm d\lambda_1(\theta),\\
 b_k&=\frac1\pi\int_0^{2\pi}G(\theta)\sin(k\theta)\,
             \mathrm d\lambda_1(\theta)\qquad(k\geq1).
 \end{aligned}
\]
Dwukrotne całkowanie przez części, z wykorzystaniem
okresowości $G$, $G'$ oraz funkcji trygonometrycznych,
daje
\[
 \begin{aligned}
 a_k&=-\frac1{\pi k^2}\int_0^{2\pi}G''(\theta)
                \cos(k\theta)\,\mathrm d\lambda_1(\theta),\\
 b_k&=-\frac1{\pi k^2}\int_0^{2\pi}G''(\theta)
                \sin(k\theta)\,\mathrm d\lambda_1(\theta).
 \end{aligned}
\]
Stąd $|a_k|+|b_k|\leq4\|G''\|_\infty/k^2$.
W szczególności szereg współczynników bezwzględnych
jest zbieżny. Poniższy argument z jądrem Poissona
pokaże również, że szereg trygonometryczny na brzegu
przedstawia dokładnie funkcję $G$, bez konieczności
przywoływania osobnego twierdzenia o zbieżności
szeregów Fouriera.

\begin{stwierdzenie}[Rozwinięcie harmoniczne na dysku]
\label{prop:pde3-fourier-disk}
Niech $g\in C^2(S^1)$ będzie rzeczywista, a
$a_0,a_k,b_k$ jej współczynnikami Fouriera z powyższą
normalizacją. Jedynym rozwiązaniem zagadnienia
$\Delta u=0$ w $B_1(0)\subset\R^2$ oraz $u|_{S^1}=g$
w klasie $C^2(B_1)\cap C(\overline{B_1})$
jest funkcja
\begin{equation}\label{eq:pde3-fourier-disk}
 u(r\cos\theta,r\sin\theta)=a_0+
 \sum_{k=1}^\infty r^k
 (a_k\cos k\theta+b_k\sin k\theta),\quad 0\leq r<1.
\end{equation}
Szereg jest bezwzględnie i jednostajnie zbieżny
na domkniętym dysku, a wszystkie jego pochodne
współrzędnych kartezjańskich można obliczać wyraz
po wyrazie na zwartych podzbiorach dysku.
\end{stwierdzenie}
\begin{proof}
Dla $0\leq r<1$ i $\psi\in\R$ szereg geometryczny
o ilorazie $re^{i\psi}$ daje
\[
 1+2\sum_{k=1}^\infty r^k\cos(k\psi)
 =\operatorname{Re}\frac{1+re^{i\psi}}{1-re^{i\psi}}
 =\frac{1-r^2}{1-2r\cos\psi+r^2}.
\]
Ostatnią równość otrzymujemy, mnożąc licznik
i mianownik ułamka zespolonego przez
$1-re^{-i\psi}$ i biorąc część rzeczywistą.
Dla ustalonego $r<1$ szereg jest jednostajnie
zbieżny względem $\psi$, gdyż $|\cos(k\psi)|\leq1$.
Dla $x=re^{i\theta}$, $\zeta=e^{it}$ mamy
$|x-\zeta|^2=1-2r\cos(\theta-t)+r^2$.
Ponieważ $\omega_2=2\pi$ oraz parametryzacja
$t\mapsto e^{it}$ ma jednostkową prędkość,
wzór na jądro Poissona z twierdzenia
\ref{thm:pde3-poisson} przyjmuje postać
\[
 P_1(re^{i\theta},e^{it})
 =\frac1{2\pi}\left(1+2\sum_{k=1}^\infty
                       r^k\cos(k(\theta-t))\right).
\]
Możemy całkować szereg wyraz po wyrazie względem
$t\in[0,2\pi]$ po pomnożeniu przez ograniczoną
funkcję $G(t)$. Korzystając z tożsamości
$\cos(k(\theta-t))=\cos(k\theta)\cos(kt)
+\sin(k\theta)\sin(kt)$, otrzymujemy dokładnie
\eqref{eq:pde3-fourier-disk}.
Jest to całka Poissona z danych $g$, a zatem
z przywołanego twierdzenia jest harmoniczna,
ciągła na domkniętym dysku i jedyna w podanej klasie.

Uzasadnijmy również bezpośrednio własności szeregu.
Każdy jego składnik jest wielomianem w zmiennych
$x,y$, ponieważ dla $z=x+iy$ zachodzi
$r^k\cos(k\theta)=\operatorname{Re}z^k$ oraz
$r^k\sin(k\theta)=\operatorname{Im}z^k$.
Dla nieujemnych liczb całkowitych $p,q$ i $m=p+q$
mamy dla $k\geq m$
\[
 \partial_x^p\partial_y^q z^k
 =i^q k(k-1)\cdots(k-m+1)z^{k-m},
\]
przy czym dla $m=0$ iloczyn jest równy $1$,
a dla $k<m$ pochodna jest zerowa.
W szczególności dla $m=2$ suma drugich pochodnych
względem $x$ i $y$ znika, co dowodzi harmoniczności
każdego składnika, także w środku dysku.

Dla $0<\rho<1$, $|z|\leq\rho$ i $k\geq\max(1,m)$
moduł pochodnej rzędu $m$ składnika szeregu
jest nie większy niż
$(|a_k|+|b_k|)k^m\rho^{k-m}$.
Szereg tych majorant jest zbieżny, ponieważ
$\sum_{k\geq\max(1,m)}k^{m-2}\rho^{k-m}<\infty$.
Z kryterium Weierstrassa otrzymujemy jednostajną
zbieżność szeregu każdej pochodnej na
$\overline{B_\rho}$. Twierdzenie o różniczkowaniu
szeregów funkcyjnych, stosowane kolejno do
pochodnych na mniejszych otwartych kulach, daje
regularność $C^\infty$ oraz możliwość obliczenia
laplasjanu wyraz po wyrazie w całym dysku.
Argument obejmuje również punkt $z=0$.

Na domkniętym dysku sam szereg jest jednostajnie
i bezwzględnie zbieżny dzięki oszacowaniu
$\sum_k(|a_k|+|b_k|)<\infty$.
Jego suma jest więc ciągła aż do brzegu.
Dla każdego $\theta$ przechodząc z $r$ do $1$
od dołu, otrzymujemy
\[
 G(\theta)=a_0+\sum_{k=1}^\infty
             (a_k\cos k\theta+b_k\sin k\theta),
\]
ponieważ granica całki Poissona jest równa
$G(\theta)$, a jednostajna zbieżność pozwala
przejść do granicy w szeregu. W szczególności
potwierdziliśmy, że suma szeregu na brzegu
jest dokładnie zadaną funkcją, a nie tylko
funkcją o tych samych formalnych współczynnikach.
\end{proof}

Współczynnik wysokiej częstości $k$ zostaje
na okręgu o promieniu $r$ pomnożony przez $r^k$.
Dla $0<r<1$ mamy $r^k=e^{-k|\log r|}$, co daje
wykładnicze tłumienie wysokich częstości.
Można zatem precyzyjnie opisać niestabilność
odtwarzania danych brzegowych.

Ustalmy $0<r<1$ i niech $T_r$ przypisuje ciągłym
danym $g$ profil $\theta\mapsto u(re^{i\theta})$,
gdzie $u$ jest całką Poissona. Jest to liniowe
odwzorowanie przestrzeni $C(S^1)$ w $C(S^1)$,
w której normą jest supremum modułu.
Jest ono różnowartościowe: jeżeli profil jest
zerowy, słaba zasada maksimum daje $u=0$
na $B_r$, a analityczność i zasada identyczności
z dowodu twierdzenia~\ref{thm:pde3-strong-max}
dają $u=0$ na $B_1$ i w konsekwencji $g=0$.

Dla $g_k(e^{i\theta})=\cos(k\theta)$, $k\geq1$,
mamy $\|g_k\|_\infty=1$ i na mocy
\eqref{eq:pde3-fourier-disk}
$T_rg_k(e^{i\theta})=r^k\cos(k\theta)$.
Zatem $\|T_rg_k\|_\infty=r^k\to0$,
podczas gdy $\|g_k\|_\infty=1$.
Odwrotność $T_r^{-1}$ określona na obrazie $T_r$
nie jest więc ciągła w zerze względem normy
supremum, nawet dla ciągu gładkich danych.
To odpowiednik niestabilności Hadamarda
z rozdziału~\ref{chap:pde-intro}.

\begin{center}
\begin{tikzpicture}[>=Stealth,scale=1]
  \draw[->] (0,0)--(5.6,0) node[right] {$k$};
  \draw[->] (0,0)--(0,3.3) node[above] {$r^k$};
  \draw[densely dotted] (0,2.9)--(5.25,2.9);
  \node[left] at (0,2.9) {$1$};
  \draw[blue!70!black,thick,domain=0:5.1,samples=100]
    plot (\x,{2.9*exp(-.25*\x)});
  \draw[red!70!black,thick,domain=0:5.1,samples=100]
    plot (\x,{2.9*exp(-.9*\x)});
  \node[blue!70!black,anchor=west] at (3.5,1.5) {$r=e^{-1/4}$};
  \node[red!70!black,anchor=west] at (3.05,.70) {$r=e^{-9/10}$};
  \node[below] at (2.8,0) {częstość $k$};
\end{tikzpicture}
\par\smallskip\small
Tłumienie modów Fouriera w rozwiązaniu harmonicznym
na dysku: im dalej od brzegu, tym szybsze zanikanie
wysokich częstości.
\end{center}

\section{Zagadnienia brzegowe na bardziej ogólnych obszarach}
\label{sec:pde3-general-domain}

Jawny wzór Poissona jest szczególnym przywilejem kuli.
Dla ogólnego obszaru funkcja Greena może nie mieć
prostej postaci, a istnienie ciągłego rozwiązania Dirichleta
zależy od geometrii brzegu. Aby rozdzielić twierdzenia
udowodnione od materiału wymagającego późniejszych narzędzi,
formułujemy poniższy wynik jako zapowiedź bez dowodu.

\begin{uwaga}[Zapowiedź twierdzenia o ogólnych obszarach]
\label{prop:pde3-future}
Niech $n\geq2$ i niech $\Omega\subset\R^n$ będzie obszarem ograniczonym o brzegu klasy $C^2$.
Wtedy dla każdego $g\in C(\partial\Omega)$ istnieje
dokładnie jedna funkcja
$u\in C^2(\Omega)\cap C(\overline\Omega)$ spełniająca
$\Delta u=0$ w $\Omega$ i $u|_{\partial\Omega}=g$.
Jest to klasyczny wynik teorii potencjału. Nie traktujemy go
tu jako wyniku udowodnionego: dowód wymaga metody Perrona
i barier brzegowych lub odpowiednich twierdzeń o regularności
eliptycznej, które wprowadzimy w dalszej części skryptu.
\end{uwaga}

Warto odróżnić dwie kwestie. \emph{Jednoznaczność}
na każdym obszarze ograniczonym wynika z zasady maksimum,
niezależnie od regularności jej brzegu. Natomiast
\emph{istnienie} dla wszystkich ciągłych danych brzegowych
nie jest prawdziwe na każdym obszarze: punkt brzegowy
może być nieregularny w sensie teorii potencjału.
To zapowiedź kryterium Wienera i pojęcia pojemności,
które pojawi się w dalszej części skryptu.

\section{Podsumowanie i mapa dalszych związków}
\label{sec:pde3-summary}

Poznaliśmy dwa wzajemnie uzupełniające się spojrzenia
na równanie Laplace'a. Własność wartości średniej,
zasady maksimum, nierówność Harnacka i analityczność
ujawniają \emph{sztywność} funkcji harmonicznych.
Z kolei rozwiązanie fundamentalne, potencjał Newtona,
funkcja Greena i jądro Poissona dają \emph{konstrukcje}
rozwiązań, a nie tylko ich jakościowe własności.

Zasadnicze pytanie pozostaje otwarte w pełnej ogólności:
jak rozwiązywać $-\Delta u=f$ na obszarze, dla której
nie znamy jawnego jądra? Odpowiedź dadzą później
przestrzenie Sobolewa, metoda wariacyjna i teoria
regularności eliptycznej. W następnym rozdziale
porównamy te własności z równaniem ciepła i jego
charakterystycznym mechanizmem wygładzania.

% Rozdział 4: równanie ciepła.
